\documentclass[11pt,a4paper]{article}
\usepackage[margin=25mm,headheight=15pt]{geometry}
\usepackage[T1]{fontenc}
\usepackage[utf8]{inputenc}
\usepackage{lmodern,microtype}
\usepackage{amsmath,amssymb,amsthm,mathtools}
\usepackage{booktabs,longtable,array,enumitem,fancyvrb}
\usepackage{xcolor,fancyhdr}
\usepackage[colorlinks=true,linkcolor=blue!35!black,citecolor=blue!35!black,urlcolor=blue!45!black,bookmarksdepth=2]{hyperref}
\newtheorem{theorem}{Theorem}[section]
\newtheorem{proposition}[theorem]{Proposition}
\newtheorem{lemma}[theorem]{Lemma}
\newtheorem{corollary}[theorem]{Corollary}
\theoremstyle{definition}
\newtheorem{definition}[theorem]{Definition}
\newtheorem{example}[theorem]{Example}
\theoremstyle{remark}
\newtheorem{remark}[theorem]{Remark}
\numberwithin{equation}{section}
\newcommand{\Qbar}{\overline{\mathbb Q}}
\newcommand{\Qb}{\overline{\mathbb Q}}
\newcommand{\QQ}{\mathbb Q}
\newcommand{\RR}{\mathbb R}
\newcommand{\CC}{\mathbb C}
\newcommand{\ZZ}{\mathbb Z}
\newcommand{\Ga}{\mathbb G_a}
\newcommand{\Gm}{\mathbb G_m}
\DeclareMathOperator{\End}{End}
\DeclareMathOperator{\Hom}{Hom}
\DeclareMathOperator{\Lie}{Lie}
\DeclareMathOperator{\Jac}{Jac}
\DeclareMathOperator{\Pic}{Pic}
\DeclareMathOperator{\Div}{Div}
\DeclareMathOperator{\Res}{Res}
\DeclareMathOperator{\Ein}{Ein}
\DeclareMathOperator{\trdeg}{trdeg}
\newcolumntype{L}[1]{>{\raggedright\arraybackslash}p{#1}}
\setlength{\emergencystretch}{3em}
\setlength{\parskip}{3pt}
\setlength{\parindent}{1.3em}
\allowdisplaybreaks[2]
\setcounter{tocdepth}{1}
\pagestyle{fancy}\fancyhf{}
\fancyhead[L]{\small Algebraic arc lengths and inverse periods}
\fancyhead[R]{\small KEIO}
\fancyfoot[C]{\thepage}
\hypersetup{pdftitle={Transcendence of Algebraic Arc Lengths and Inverse Periods},pdfauthor={KEIO},pdfsubject={Division dichotomies, beta quantiles, and arithmetic transfers}}
\title{\bfseries Transcendence of Algebraic Arc Lengths\\and Inverse Periods\\[6pt]
\large Division Dichotomies, Beta Quantiles,\\and Arithmetic Transfers}
\author{KEIO}
\date{12 October 2026}
\begin{document}
\maketitle
\begin{abstract}
We develop arithmetic consequences of algebraic-integral transcendence,
consolidating the research discussion subsequent to
\emph{Explicit Transcendence Criteria for Beta-Type Integrals}.
For noncircular polynomial lemniscates, the square root of the complexified
Euclidean metric becomes a holomorphic differential on a finite cover.
This yields transcendence and logarithmic separation for closed lengths
and arcs with algebraic endpoints. We classify rational division parameters
for all positive rational sinusoidal-spiral exponents, equivalently all
rational quantiles of the beta distributions with parameters
$(\alpha,1/2)$, $\alpha\in\mathbb Q_{>0}$. Torsion-packet results strengthen
selected cases to all algebraic division ratios. For a general positive
regular differential, a minimal abelian quotient separates elliptic
division from a finite set of algebraic exceptions; a division-hull
argument treats algebraic, rather than merely rational, ratios.
Explicit Cassini and hyperelliptic constructions give infinite families
linearly independent over $\Qbar$. Further transfers concern quadrature,
orthogonal polynomials, probability, differential equations, and spectral
models. Established transcendence and arithmetic-geometric theorems are
identified as inputs. Linear independence is distinguished from algebraic
independence, and unresolved constant problems remain unresolved.
\end{abstract}
\noindent\textbf{Keywords.} Abelian integrals; generalized Jacobians;
lemniscates; beta quantiles; transcendence; analytic subgroup theorem;
torsion points; quadrature.

\section*{Scope, provenance, and conventions}
\addcontentsline{toc}{section}{Scope, provenance, and conventions}
This sequel covers the discussion after the manuscript
\path{Transcendence_Beta_Gamma_Unified.pdf} (11 October 2026).
Necessary inputs are restated, and overlapping applications are grouped
under common theorems. The numbered observations are not a count of
independent research breakthroughs.

This is a research manuscript; neither external peer review nor
historical priority is claimed. The analytic subgroup theorem,
logarithmic and $E$-function theorems, torsion classifications, and
arithmetic finiteness theorems are established inputs, identified in the
proofs. Neither Euler's constant nor an unresolved individual Gamma
value such as $\Gamma(1/5)$ is proved transcendental here.

We fix $\Qbar\subset\mathbb C$. Real algebraic data lie in
$\Qbar\cap\mathbb R$, and positive branches are used on real integration
intervals. Curves mean smooth projective models unless specified;
``regular'' means globally holomorphic. An oval is a real connected
component, not necessarily an elliptic curve. Traversals at singular
plane points are interpreted on the normalization. Linear independence
means independence over $\Qbar$, and an infinite family is independent
when every finite subfamily is. This differs from algebraic independence.
Algebraicity of a point means algebraicity of all its coordinates;
coordinate-wise transcendence requires its own proof.

\clearpage
\begingroup\small\setlength{\parskip}{0pt}\tableofcontents\endgroup
\clearpage

\part{Algebraic lengths and their arithmetic}

\section{Polynomial lemniscates: holomorphic length differentials}
\label{lm-section}

The baseline argument identifies the length of a polynomial lemniscate with
a nonzero closed period.  The additional observation in this section is
stronger: except for circles, the length differential can be made
\emph{holomorphic on a complete algebraic curve}.  This removes the
closed-path restriction, separates the resulting lengths from logarithms,
and permits an explicit infinite linear-independence construction.
Tamiozzo's Remark~8.16 asks whether all polynomial lemniscates over
$\Qbar$ have transcendental length; the examples proved there include
$|z^k-a^k|=1$ \cite{lm-Tamiozzo}.  We give the geometric construction
explicitly, including critical levels.  The transcendence inputs throughout
are established forms of the analytic subgroup theorem, not new versions of
that theorem.

Write $\mathbb R_{\rm alg}=\Qbar\cap\mathbb R$.  Let
\begin{equation}\label{lm-level}
 P\in\Qbar[z]\setminus\Qbar,\qquad
 \rho\in\mathbb R_{\rm alg},\quad \rho>0,\qquad
 \Lambda(P,\rho)=\{z\in\mathbb C:|P(z)|=\rho\}.
\end{equation}
Length always means ordinary Euclidean geometric length, with overlapping
arcs counted once.  A finite arc family below means finitely many arcs
whose interiors are disjoint up to finitely many points.  An algebraic
endpoint has both Cartesian coordinates in $\mathbb R_{\rm alg}$.

\subsection{Normalization and the positive square-root lift}

After replacing $P$ by $P/\rho$, assume that $\rho=1$.  If $d$ is the
greatest common divisor of the multiplicities of the distinct roots, choose
$Q\in\Qbar[z]$ with $P=Q^d$.  Then $\Lambda(P,1)=\Lambda(Q,1)$ and the
root multiplicities $m_1,\ldots,m_r$ of $Q$ have greatest common divisor
one.  Put $n=\deg Q$ and let $Q^\#(w)$ denote the polynomial obtained by
conjugating the coefficients of $Q$.

\begin{lemma}[The underlying curve]\label{lm-normalization}
The affine complex curve
\begin{equation}\label{lm-complexification}
 Q(z)Q^\#(w)=1
\end{equation}
is geometrically irreducible.  Its smooth projective normalization $X$
has a real structure under $z=x+iy$, $w=x-iy$.  Its real locus is a finite
union of circles, all lying over finite points of the affine plane.  On
each such circle, the map to the plane is an immersion, including over
critical points of $Q$ on the lemniscate.
\end{lemma}

\begin{proof}
Work over $\Qbar(z)$ and consider the degree-$n$ polynomial
$Q^\#(w)-1/Q(z)$, scaled to be monic.  At a zero $z_j$ of multiplicity
$m_j$, every root $\xi$ has valuation $-m_j/n$: its valuation must be
negative, and the leading term of $Q^\#(\xi)$ is then uniquely of least
valuation.  If a monic factor has degree $e$, its nonzero constant term
therefore has valuation $-em_j/n$.  This is an integer, so $n$ divides
$em_j$ for every $j$.  Since $\gcd_j m_j=1$, $n$ divides $e$.  A proper
factor is impossible.

The leading homogeneous part of the real plane equation is
$|a_n|^2(x^2+y^2)^n$, so its projective closure has no real point at
infinity.  Consequently the real normalization is compact and its
components are circles.  To examine a point $z_0$ on the level set, take a
local holomorphic logarithm of $Q(z)/Q(z_0)$.  It vanishes to some order
$m\ge1$ and, after a biholomorphic change of local coordinate, equals
$u^m$.  The real level set is thus
$\operatorname{Re}(u^m)=0$, a union of $m$ distinct real lines.  The
normalization separates these branches, and the inverse coordinate change
has nonzero derivative on each.  Hence the normalized map is an immersion.
This also shows that every affine real point of the lemniscate is covered
by real branches of the normalization.
\end{proof}

On $X$ consider the rational quadratic differential
\begin{equation}\label{lm-quadratic}
 q=dz\,dw=dx^2+dy^2.
\end{equation}
Choose a nonzero rational differential $\tau$ on $X$, write
$q=f\tau^2$, and normalize the quadratic extension obtained by adjoining
$\sqrt f$.  If $f$ is already a square, use $X$ itself and choose a square
root.  We obtain a connected smooth projective curve $Y$ over $\Qbar$, a
map $\pi:Y\to X$ of degree at most two, and a rational differential
$\omega$ satisfying
\begin{equation}\label{lm-square}
 \omega^2=\pi^*q.
\end{equation}
All choices here are algebraic.  On the real normalization circles,
$q$ is strictly positive on nonzero real tangent vectors.  Thus the cover
is unramified there, and, after orienting a circle, exactly one local lift
has $\omega$ positive on its tangent.  These local lifts join to a closed
lift: the positive choice is unique when the initial point and tangent
return.  In the degree-one case one chooses the orientation giving a
positive integral.  In either case its integral is its Euclidean length.

\begin{proposition}[Closed lengths]\label{lm-closed-lengths}
The length of every geometric connected component of
$\Lambda(P,\rho)$, and its total length, is a nonzero period of a rational
differential on an algebraic curve over $\Qbar$.  Consequently each is
transcendental.
\end{proposition}

\begin{proof}
The normalization is one-to-one away from finitely many singular points.
For a geometric connected component, take all the normalization circles
lying above it, lift them as just described, and add their cycles.  Their
integrals give the geometric length, since the exceptional finite set has
length zero.  This integral is positive.  W\"ustholz's theorem that a
nonzero closed period of an algebraic rational differential is
transcendental applies; see \cite[Theorem~8.14]{lm-Tamiozzo} and
\cite{lm-HuberWustholz}.  Repeated factors were removed without changing
the geometric set, so no multiplicity is counted.
\end{proof}

\subsection{The additional holomorphy statement}

\begin{lemma}[Holomorphy away from the circle case]\label{lm-holomorphy}
If $P$ has at least two distinct roots, the differential $\omega$ in
\eqref{lm-square} is holomorphic everywhere on $Y$.
\end{lemma}

\begin{proof}
At finite points of $X$, both $z$ and $w$ are regular, so $q=dz\,dw$ has
no pole.  To inspect infinity, use the closure of
\eqref{lm-complexification} in $\mathbb P^1\times\mathbb P^1$.  If $z$
has a pole at a point of $X$, the defining equation forces $w$ to tend to
a finite root $\beta$ of $Q^\#$.  Write
\[
 \operatorname{ord}(z)=-a,\qquad
 \operatorname{ord}(w-\beta)=b,\qquad a,b>0,
\]
and let $m$ be the multiplicity of $\beta$.  Comparing valuations gives
$na=mb$.  Because $Q$ has at least two distinct roots, $m<n$, whence
$b>a$.  In characteristic zero,
\begin{equation}\label{lm-order}
 \operatorname{ord}(q)=(-a-1)+(b-1)=b-a-2\ge-1.
\end{equation}
The case where $w$ has a pole is identical.  These exhaust the points at
infinity; the two coordinates cannot both have poles.

If a quadratic differential has local order $k\ge-1$, an even $k$ gives
an unramified square root of order $k/2\ge0$.  For odd $k$, the square-root
cover has local parameter $t=s^2$, and the pulled-back quadratic
differential has order $2k+2$.  Its square root has order $k+1\ge0$.
Thus even a possible simple pole of $q$ produces a regular $\omega$.
\end{proof}

The calculation at infinity is essential.  The fact that a length is a
closed period alone would not justify replacing its cycle by an open arc.
Lemma~\ref{lm-holomorphy} supplies precisely the stronger input needed for
that replacement.

\subsection{Separation of Abelian integrals from logarithms}

Define the $\Qbar$-vector space
\begin{equation}\label{lm-logspace}
 \mathcal L=\Qbar+
 \operatorname{span}_{\Qbar}
 \{\lambda\in\mathbb C:\exp(\lambda)\in\Qbar^\times\}.
\end{equation}
All branches of logarithms are allowed; in particular $2\pi i\in
\mathcal L$.  A first-kind algebraic integral means an integral of a
holomorphic differential defined over $\Qbar$ on a smooth projective
curve over $\Qbar$, along a path with algebraic endpoints.  Closed
periods are included.  Finite sums can equivalently be treated on a
product of Jacobians.

\begin{lemma}[Abelian--logarithmic separation]\label{lm-separation}
Let $S$ be a $\Qbar$-linear combination of first-kind algebraic integrals.
If $S\in\mathcal L$, then $S=0$.  Thus every nonzero such $S$ is
transcendental and $S/\pi$ is transcendental.
\end{lemma}

\begin{proof}
On a product $A$ of the relevant Jacobians, write $S=\eta(v)$, where
$\eta$ is an invariant algebraic differential and
$\exp_A(v)\in A(\Qbar)$.  This uses the algebraicity of the endpoint
divisors.  Suppose
\[
 \eta(v)+c_0+\sum_{j=1}^r c_j\lambda_j=0,
 \qquad c_j\in\Qbar,\quad e^{\lambda_j}\in\Qbar^\times.
\]
Apply the refined analytic subgroup theorem
\cite[Theorem~6.2]{lm-HuberWustholz} to
\[
 G=A\times\mathbb G_a\times\mathbb G_m^r,
 \qquad u=(v,1,\lambda_1,\ldots,\lambda_r).
\]
Its exponential is algebraic.  The theorem gives a connected algebraic
subgroup $H\le G$ with $u\in\operatorname{Lie}(H)_{\mathbb C}$ and with
the displayed algebraic linear form vanishing on $\operatorname{Lie}(H)$.

A connected subgroup of a product of an abelian variety and a connected
affine commutative algebraic group splits between these two types.  For
example, the algebraic-group Goursat argument reduces a possible mixing
to a common connected quotient, which would be both proper and affine
and hence trivial.  Accordingly $H=B\times H_{\rm aff}$ for an abelian
subvariety $B\le A$.  Restricting the linear form to $\operatorname{Lie}B$
shows that $\eta$ vanishes there.  Since $v\in\operatorname{Lie}(B)$,
$S=0$.  The last conclusions follow because both $\Qbar$ and
$\Qbar\pi$ lie in $\mathcal L$.
\end{proof}

\begin{theorem}[Algebraic-endpoint arcs and finite arc families]
\label{lm-arcs}
Suppose that $P$ has at least two distinct roots.  The length $L>0$ of any
finite arc family on $\Lambda(P,\rho)$ with algebraic endpoints is a
finite sum of first-kind algebraic integrals.  Consequently
\begin{equation}\label{lm-arcs-conclusion}
 L\notin\mathcal L,
 \qquad L\text{ is transcendental},\qquad
 L/\pi\text{ is transcendental}.
\end{equation}
The same conclusions hold for each full geometric connected component
and for the total length.
\end{theorem}

\begin{proof}
Split an arc, when necessary, at its finitely many critical crossings.
Those crossing points are algebraic: their complex coordinate is a zero
of $P'$.  Normalize each resulting branch and choose the lift for which
$\omega$ integrates positively.  Every lift of an algebraic endpoint
under these algebraic finite maps is algebraic.  The resulting integrals
are of the first kind by Lemma~\ref{lm-holomorphy}, and their sum is $L$.
For complete components, the earlier closed positive lifts give the same
conclusion directly.  Positivity and Lemma~\ref{lm-separation} finish the
proof.
\end{proof}

\begin{corollary}[Individual lobes and boundary circuits]
\label{lm-lobes}
The conclusions of Theorem~\ref{lm-arcs} apply to each positive-length
lobe or finite boundary circuit of a polynomial lemniscate whose vertices
are critical crossings.  In particular, each bounded complementary
region has transcendental boundary length.
\end{corollary}

\begin{proof}
At a crossing, a planar circuit may switch normalized branches; it need
not lift to a closed cycle.  Instead, cut at its algebraic vertices and
apply the finite-arc assertion.  Boundary arcs are counted once when
computing geometric boundary length.  A boundary without vertices is a
complete smooth component and was already covered.
\end{proof}

For a circle, positive arcs with algebraic endpoints are also
transcendental in length, although \eqref{lm-arcs-conclusion} need not
hold in full.  Indeed the radius is algebraic and an arc length is its
product with an angle $\theta$ such that $e^{i\theta}$ is algebraic.
A nonzero algebraic linear combination of logarithms of algebraic
numbers cannot be algebraic, by the additive--multiplicative case of the
same subgroup theorem.  Thus the transcendence assertion for algebraic
endpoint arcs holds for \emph{every} polynomial in \eqref{lm-level}.

\begin{theorem}[Exactly when a perimeter is an algebraic multiple of $\pi$]
\label{lm-circle-dichotomy}
Let $L$ be the total length or the length of any geometric connected
component of $\Lambda(P,\rho)$.  Then
\[
 L/\pi\in\Qbar
 \quad\Longleftrightarrow\quad
 P(z)=a(z-b)^n\text{ for some }a\in\Qbar^\times,
 \ b\in\Qbar,\ n\ge1.
\]
Equivalently, this is precisely the circle case.
\end{theorem}

\begin{proof}
If $P=a(z-b)^n$, the level is the circle of radius
$R=(\rho/|a|)^{1/n}\in\mathbb R_{\rm alg}$, so $L=2\pi R$.
Otherwise $P$ has at least two distinct roots and
Theorem~\ref{lm-arcs} excludes algebraicity of $L/\pi$.

For completeness, the geometric description of the exception is also
elementary.  If a circle $z=b+R\zeta$, $|\zeta|=1$, is contained in the
level set, put $F(\zeta)=P(b+R\zeta)$ and $N=\deg F$.  Define
$F^*(\zeta)=\zeta^N\overline{F(1/\bar\zeta)}$.
The circle identity implies the polynomial identity
$F(\zeta)F^*(\zeta)=\rho^2\zeta^N$.  Hence $F$ has no nonzero root,
so $F(\zeta)=c\zeta^N$.  Thus $P$ has the asserted single-root form.
This argument also excludes a circle component in a noncircle case.
\end{proof}

\begin{corollary}[Exponentials and trigonometric values]\label{lm-exponentials}
Let $L>0$ be any length covered by Theorem~\ref{lm-arcs}, and let
$c\in\Qbar^\times$.  Then
\[
 e^{cL},\qquad \sin(cL),\qquad \cos(cL),\qquad \tan(cL)
\]
are transcendental; the last expression is defined.
\end{corollary}

\begin{proof}
If $e^{cL}$ were algebraic, $cL$ would belong to the logarithmic space
\eqref{lm-logspace}, contrary to Lemma~\ref{lm-separation}.  An algebraic
sine or cosine would make $e^{icL}$ a root of a quadratic polynomial over
$\Qbar$.  Thus both are transcendental and in particular $\cos(cL)\ne0$.
Finally an algebraic tangent would make
$e^{2icL}=(1+i\tan(cL))/(1-i\tan(cL))$ algebraic.  The denominator cannot
vanish for a finite exponential representation.  Each alternative is
impossible.
\end{proof}

\begin{theorem}[Elementary integration occurs exactly for circles]
\label{lm-elementary}
The algebraic length differential of a polynomial lemniscate admits an
elementary antiderivative, in the differential-algebraic sense, if and
only if the lemniscate is a circle.  Elementary extensions and algebraic
changes of variable are allowed; a numerical special value of an
integral is not what is being classified here.
\end{theorem}

\begin{proof}
For a circle of algebraic radius $R$ the length differential is, up to
orientation, $-iR\,d\zeta/\zeta$, and its primitive is logarithmic.
Otherwise Lemma~\ref{lm-holomorphy} gives a nonzero holomorphic
differential on a complete smooth curve.  Such a differential has no
elementary primitive.  Indeed Liouville's theorem \cite{lm-Rosenlicht}
would give, on a finite algebraic cover if necessary, an identity
\[
 \omega=dF+\sum_{j=1}^r c_j\frac{dG_j}{G_j},
 \qquad c_j\in\mathbb C,
\]
with $F,G_j$ meromorphic.  A pole of $F$ gives a pole of order at least
two in $dF$, whereas logarithmic derivatives have only simple poles.
Thus $F$ is constant.  Vanishing of the remaining residues means
$\sum_j c_j\operatorname{div}(G_j)=0$.  This system has integer
coefficients, so its complex kernel is spanned by rational relations.
Clearing denominators in each rational relation makes the corresponding
product of powers of the $G_j$ constant, so its logarithmic derivative
is zero.  Consequently $\omega=0$, a contradiction.  Pullback to a
finite cover preserves holomorphy and nonvanishing.
\end{proof}

\begin{proposition}[Linear combinations of lengths]\label{lm-linear-combinations}
For finitely many polynomial lemniscate lengths, including the
algebraic-endpoint arc lengths above, every $\Qbar$-linear combination is
either zero or transcendental.  For lengths drawn only from polynomials
with at least two distinct roots, a nonzero combination lies outside
$\mathcal L$ and its quotient by $\pi$ is transcendental.
\end{proposition}

\begin{proof}
Separate the combination into first-kind integrals from noncircle
lemniscates and logarithmic terms from circles.  If the total were
algebraic, Lemma~\ref{lm-separation} would force the first-kind part to
vanish.  The remaining logarithmic linear form is either zero or
transcendental, so an algebraic total must be zero.  The second assertion
is exactly Lemma~\ref{lm-separation} on a product of Jacobians.
\end{proof}

These statements do not assert algebraic independence of lengths, or even
linear independence of an arbitrary chosen collection.  Isometries,
algebraic dilations and symmetries create genuine dependencies.  The next
construction provides an explicit collection for which independence can
be proved.

\section{Closed-period transfers to other geometric quantities}
\label{lm-geometric-transfers}

The following extensions use rational differentials and closed cycles;
they do not require the extra holomorphy for noncircle polynomial
lemniscates.  A finite algebraic linear combination of closed
rational-differential periods, even on different curves, is either zero
or transcendental.  To justify this without assuming that sums of
transcendental numbers are transcendental, puncture each curve at the
poles, take their product, and use the closed algebraic $1$-form
$\sum_j c_j\operatorname{pr}_j^*\eta_j$ there.  Its period on the
corresponding homology class is the required sum.  W\"ustholz's theorem
for smooth quasi-projective varieties applies
\cite[Theorem~3.1]{lm-Braden}.

\begin{theorem}[Rational lemniscates]\label{lm-rational}
Let $R\in\Qbar(z)$ be nonconstant and $\rho>0$ algebraic.  Every bounded
geometric connected component of
$\{z\in\mathbb C:|R(z)|=\rho\}$ has transcendental Euclidean length.
On the Riemann sphere, every connected component and the whole level set
have transcendental length for the unit-sphere metric
\[
 ds_{\rm sph}=\frac{2|dz|}{1+|z|^2}.
\]
The spherical level includes $\infty$ when $R(\infty)$ is finite and
has modulus $\rho$.
\end{theorem}

\begin{proof}
Normalize the reduced algebraic correspondence
$R(z)R^\#(w)=\rho^2$, component by component, taking its closure from the
locus where both rational functions are defined and finite.  Any irreducible component
supporting a real arc is fixed by the real structure.  At each level
point, a local logarithm of $R/R(z_0)$ reduces the real equation to
$\operatorname{Re}(u^m)=0$.  The normalized branches are therefore
immersed, including at critical crossings and, in a spherical local
coordinate, at infinity.

For a bounded Euclidean component the normalization circles remain in
the finite plane; use $q=dz\,dw$.  For spherical length use
\[
 q_{\rm sph}=\frac{4\,dz\,dw}{(1+zw)^2}.
\]
This is positive and nonsingular along the real locus.  At infinity,
$u=1/z$, $v=1/w$ give the identical expression
$4\,du\,dv/(1+uv)^2$.  Each positive square-root lift is closed, giving
a rational-differential period.  Sums over different normalization
components are covered by the product argument above.  Positivity
proves transcendence.  No finite Euclidean length is asserted for an
unbounded component.
\end{proof}

For the next two results, an algebraic immersed closed curve means a
real normalization circle $\gamma$ on a smooth projective curve over
$\mathbb R_{\rm alg}$, with rational coordinate functions $x,y$ regular
along $\gamma$, such that $(dx,dy)$ never vanishes there.  Integrals
count one traversal.  For an embedding, or a map generically one-to-one
onto its image, this is geometric length without repeated arcs.

\begin{theorem}[Weighted length and bending energies]\label{lm-weighted}
For such a curve, let $W\in\mathbb R_{\rm alg}[x,y]$ be nonnegative on
the image and not identically zero there.  Then
\[
 \int_\gamma W\,ds
 \quad\text{and}\quad
 \int_\gamma\kappa^{2m}\,ds\quad(m\ge1)
\]
are transcendental, where $\kappa$ is Euclidean curvature.
\end{theorem}

\begin{proof}
Use the positive closed lift with algebraic differential $\omega=ds$.
The weighted length is the positive period of the rational differential
$W\omega$.  For curvature, the derivation $D$ defined by
$df=(Df)\omega$ is algebraic on the square-root curve and is
arc-length differentiation on the real lift.  Hence
\[
 \kappa^2=(D^2x)^2+(D^2y)^2
\]
is rational and regular along that lift.  The bending energy is the
period of $(\kappa^2)^m\omega$.  It is positive, because identically
zero curvature would make the unit tangent constant, which is
incompatible with a nonconstant closed immersed curve.  The
closed-period theorem proves both assertions.
\end{proof}

\begin{theorem}[Anisotropic length derivatives]\label{lm-anisotropic}
For an algebraic immersed closed curve define
\[
 A(t)=\int_\gamma\sqrt{t\,dx^2+dy^2},\qquad t>0.
\]
If $x$ is nonconstant, then for every positive algebraic $t$ and every
integer $n\ge0$, $A^{(n)}(t)$ is nonzero and transcendental.  In this
planar closed-immersion setting, nonconstancy of $x$ is automatic.
\end{theorem}

\begin{proof}
Positivity of the metric on the compact immersed curve permits
differentiation under the integral to every order.  Fix an algebraic
$t>0$ and its positive closed square-root differential
$\omega_t^2=t\,dx^2+dy^2$.  For $n\ge1$,
\begin{equation}\label{lm-anisotropic-formula}
 A^{(n)}(t)=c_n\int_{\widetilde\gamma}
       \left(\frac{dx}{\omega_t}\right)^{2n}\omega_t,
 \qquad c_n=\prod_{j=0}^{n-1}\left(\frac12-j\right)\in\mathbb Q^\times.
\end{equation}
The integrand is rational over $\Qbar$ and pole-free on the cycle.  The
integral before multiplication by $c_n$ is strictly positive because
$dx$ is not identically zero.  It is consequently a nonzero period, as
is the length for $n=0$.  If $x$ were constant, the periodic real
function $y$ would have a critical point, contradicting immersion.
Without that immersion assumption the coordinate hypothesis must be
retained: otherwise all positive-order derivatives can vanish.
\end{proof}

\begin{theorem}[Areas bounded by smooth algebraic ovals]\label{lm-area}
Let $\Omega\subset\mathbb R^2$ be a nonempty bounded domain whose boundary
is a finite disjoint union of smooth embedded algebraic ovals over
$\mathbb R_{\rm alg}$.  If $W\in\mathbb R_{\rm alg}[x,y]$ is
nonnegative on $\Omega$ and not identically zero there, then
\[
 \iint_\Omega W(x,y)\,dx\,dy
\]
is transcendental.  This includes the ordinary area.
\end{theorem}

\begin{proof}
The polynomial $F(x,y)=\int_0^xW(s,y)\,ds$ has algebraic coefficients.
Green's theorem with the induced boundary orientation gives
\[
 \iint_\Omega W\,dx\,dy=\int_{\partial\Omega}F\,dy.
\]
Each oval is a closed cycle on its algebraic normalization; $F\,dy$ is
rational and has no pole on that cycle.  The sum is strictly positive,
so the product form of the closed-period theorem applies.  Inner
boundaries carry the opposite orientation.  Positivity comes from the
area integral, not from each separate boundary integral.
\end{proof}

The smooth-boundary hypothesis is substantive.  The critical filled
lemniscate $\Omega_*=\{z:|z^2-1|\le1\}$ has area $2$.  Its polar
equation gives $0\le r^2\le2\cos(2\theta)$ where the cosine is
nonnegative, and hence
\[
 \operatorname{area}(\Omega_*)=
 \int_{\{0\le\theta\le2\pi:\cos(2\theta)\ge0\}}
          \cos(2\theta)\,d\theta=2.
\]
The positive boundary orientation switches branches at the node and
does not produce the required closed cycle on the normalization.  Thus
we do not assert transcendental area for every singular algebraic
region.

\section{An infinite linearly independent family of Cassini perimeters}
\label{lm-cassini-section}

For each odd prime $p$, let
\begin{equation}\label{lm-cassini-family}
 C_p=\{z\in\mathbb C:|z^2-p^{-2}|=1\},
 \qquad \ell_p=\operatorname{length}(C_p),
 \qquad m_p=\frac{1-\sqrt{1-p^{-4}}}{2}.
\end{equation}
Each $C_p$ is a smooth connected Cassini oval.  We use the parameter,
rather than the square-root modulus, convention
\[
 K(m)=\int_0^1\frac{dt}{\sqrt{(1-t^2)(1-mt^2)}}.
\]
The classical reduction of the Cassini perimeter, in the form proved in
\cite[Equation~(7.7.4)]{lm-Tamiozzo}, gives
\begin{equation}\label{lm-cassini-period}
 \ell_p=4K(m_p).
\end{equation}
Thus $\ell_p$ is a nonzero complete period of the algebraic differential
$dx/y$ on the elliptic curve
\begin{equation}\label{lm-elliptic-curve}
 E_p:\quad y^2=x(1-x)(1-m_px).
\end{equation}
Indeed $\int_0^1dx/y=2K(m_p)$, and traversing this interval on one sheet
and returning on the other doubles the integral.

\begin{lemma}[Pairwise nonisogeny]\label{lm-nonisogeny}
The elliptic curves $E_p$, as $p$ ranges over the odd primes, are pairwise
nonisogenous over $\Qbar$.
\end{lemma}

\begin{proof}
The standard Legendre invariant is
\[
 j(m)=256\frac{(1-m+m^2)^3}{m^2(1-m)^2}.
\]
Since $m_p(1-m_p)=1/(4p^4)$, direct substitution yields
\begin{equation}\label{lm-j-invariant}
 j(E_p)=64\frac{(4p^4-1)^3}{p^4}\in\mathbb Q.
\end{equation}
At the rational prime $p$ this has valuation $-4$.  For a different odd
prime $q$, the rational number $j(E_q)$ has nonnegative valuation at $p$.

An elliptic curve over a number field has potentially good reduction at
a finite place if and only if its $j$-invariant is integral there; see
\cite[Chapter~VII, Proposition~5.5]{lm-Silverman}.  Potentially good
reduction is invariant under isogeny.  One way to see the latter assertion
is to pass to an extension with good reduction and use the
N\'eron--Ogg--Shafarevich criterion: an isogeny identifies the rational
$\ell$-adic Tate modules for an auxiliary prime $\ell$ different from the
residue characteristic \cite{lm-SerreTate}.

If $E_p$ and $E_q$ were isogenous over $\Qbar$, the isogeny and both curves
would be defined over a common number field.  At any place over $p$,
$E_q$ would have potentially good reduction and $E_p$ would not.  This is
impossible.  The argument allows arbitrary degrees of isogeny and all
finite extensions of the field of definition.
\end{proof}

\begin{theorem}[Cassini perimeter independence]\label{lm-cassini-independence}
The collection
\begin{equation}\label{lm-independent-family}
 \{1,\pi\}\ \cup\ \{\ell_p:p\text{ an odd prime}\}
\end{equation}
is linearly independent over $\Qbar$.  Here an infinite collection is
independent when each of its finite subcollections is independent.
\end{theorem}

\begin{proof}
Take distinct odd primes $p_1,\ldots,p_r$.  For each $j$, choose a period
vector $v_j\in\operatorname{Lie}(E_{p_j})_{\mathbb C}$ on which the
invariant differential in \eqref{lm-elliptic-curve} takes the value
$\ell_{p_j}$.  In particular $v_j\ne0$ and $\exp(v_j)=0$.
Suppose
\[
 c_0+c_\pi\pi+\sum_{j=1}^r c_j\ell_{p_j}=0,
 \qquad c_0,c_\pi,c_j\in\Qbar.
\]
Apply the analytic subgroup theorem to
\[
 G=\prod_{j=1}^r E_{p_j}\times\mathbb G_a\times\mathbb G_m,
 \qquad u=(v_1,\ldots,v_r,1,2\pi i),
\]
and to the invariant linear form with coefficients
$c_1,\ldots,c_r,c_0,c_\pi/(2i)$.  It yields a connected subgroup $H$
containing $u$ in its Lie algebra and annihilated by that form.

The abelian and affine parts split as in
Lemma~\ref{lm-separation}.  By Lemma~\ref{lm-nonisogeny}, the elliptic
factors are simple and pairwise nonisogenous.  Hence each connected
abelian subvariety of their product is a product of some of those
factors: a nonproduct subvariety would yield a nonzero homomorphism
between two distinct simple factors, and therefore an isogeny.  Every
projection of $u$ to an elliptic Lie algebra is nonzero, so all elliptic
factors must occur in $H$.

Likewise a connected subgroup of $\mathbb G_a\times\mathbb G_m$ splits
between its additive and multiplicative factors.  Its projections must
both be nonzero since the corresponding coordinates of $u$ are $1$ and
$2\pi i$.  Thus $H=G$.  The purported linear form therefore vanishes on
all of $\operatorname{Lie}G$, forcing every coefficient to be zero.
\end{proof}

\begin{corollary}\label{lm-cassini-ratios}
For distinct odd primes $p,q$, the quotient $\ell_p/\ell_q$ is
transcendental.  For any finite choice of distinct odd primes, every
nonzero algebraic linear combination of their perimeters is
transcendental; adjoining $1$ and $\pi$ introduces no linear relation.
\end{corollary}

\begin{proof}
An algebraic quotient would give a forbidden two-term linear relation.
If a nonzero linear combination of the perimeters were algebraic, moving
that algebraic number to the constant term would contradict
Theorem~\ref{lm-cassini-independence}.
\end{proof}

This is linear independence, not algebraic independence.  No assertion
about products of distinct perimeters, their polynomial relations, or
the algebraic independence of $\pi$ and any one perimeter is needed.
The following register matches the numbered post-baseline results to
the precise statements here.  Their hypotheses are retained, especially
the noncircle restriction in entries~62--63 and the smooth-boundary
restriction in entry~68.
\begin{center}
\begin{tabular}{p{0.07\linewidth}p{0.83\linewidth}}
\hline
61 & Algebraic-endpoint arc lengths: Theorem~\ref{lm-arcs} and the
circle discussion following Corollary~\ref{lm-lobes}.\\
62 & Noncircle length divided by $\pi$: Theorems~\ref{lm-arcs} and
\ref{lm-circle-dichotomy}.\\
63 & Exponential, sine and cosine values for noncircle lengths:
Corollary~\ref{lm-exponentials}, also treating tangent.\\
64 & Elementary antiderivative exactly for circles:
Theorem~\ref{lm-elementary}.\\
65 & Rational-lemniscate Euclidean and spherical lengths:
Theorem~\ref{lm-rational}.\\
66 & Weighted lengths and even bending energies:
Theorem~\ref{lm-weighted}.\\
67 & Anisotropic length derivatives: Theorem~\ref{lm-anisotropic}.\\
68 & Areas with smooth algebraic oval boundaries:
Theorem~\ref{lm-area}.\\
69 & Cassini perimeter independence:
Theorem~\ref{lm-cassini-independence}.\\
\hline
\end{tabular}
\end{center}

% lm-DONE


\part{Division and inverse periods}

\section{Arithmetic division of real Abelian integrals}
\label{dd-section}

This section distinguishes three questions: which rational fractions of an
Abelian integral have algebraic division points, whether replacing rational
fractions by arbitrary algebraic fractions changes finiteness, and whether
the resulting genus bounds can be improved.  The answers are stated for
holomorphic differentials on real algebraic curves.  Consequently they apply
to an arc-length differential only after its holomorphic realization on a
smooth algebraic length cover has been established.  A real oval means a
connected component of the real locus of a smooth curve; it is a topological
circle, not an elliptic curve.

Put $k=\Qbar\cap\mathbb R$.  Let $C/k$ be a smooth projective geometrically
connected curve of genus $g$, and let
$0\ne\omega\in H^0(C,\Omega^1_{C/k})$.  We consider either of the following
two settings.
\begin{enumerate}
\item A real oval $\Gamma\subset C(\mathbb R)$ on which $\omega$ is nowhere
zero, oriented so that its value on the positive tangent is positive.  Fix
$P_0\in\Gamma(k)$ and traverse the oval once, returning to $P_0$.
\item An embedded compact arc $\Gamma\subset C(\mathbb R)$ with distinct
endpoints $P_0,P_1\in C(k)$, oriented from $P_0$ to $P_1$, on which $\omega$
is positive and nowhere zero, including at the endpoints.
\end{enumerate}
Let $L=\int_\Gamma\omega>0$.  For $0\le t\le1$, let $P_t$ be the unique
point reached after integral $tL$ along the chosen traversal.  In the oval
case $P_1=P_0$, and we count parameters in $I=[0,1)$; in the arc case
$I=[0,1]$.  Thus the identification at the end of a closed traversal is
never counted twice.  Define
\begin{equation}\label{dd-sets}
 S_{\mathbb Q}=\{t\in I\cap\mathbb Q:P_t\in C(k)\},\qquad
 S_{\rm alg}=\{t\in I\cap k:P_t\in C(k)\}.
\end{equation}
Positivity makes the parameter $t$ single-valued.  It also makes the
traversal real analytic, because a local primitive of the nonvanishing
differential is a real analytic coordinate.

\subsection{The minimal Abelian quotient and a logarithmic rigidity lemma}

Let $J=\operatorname{Jac}(C)$ and let $j:C\to J$ be the Abel--Jacobi map
$P\mapsto[P-P_0]$.  Write $\Omega$ for the invariant differential on $J$
whose pullback is $\omega$.

\begin{lemma}\label{dd-minimal}
There is a unique largest Abelian subvariety $B\subset J_{\Qbar}$ on
which $\Omega$ vanishes.  It is defined over $k$.  The quotient
\[
 \pi:J\longrightarrow A:=J/B
\]
and a nonzero invariant differential $\eta$ on $A$ are defined over $k$,
with $\Omega=\pi^*\eta$.  No positive-dimensional Abelian subvariety of
$A_{\Qbar}$ is annihilated by $\eta$.  Moreover,
\[
 f:=\pi\circ j:C\longrightarrow A
\]
is nonconstant, its image contains $0$ and generates $A$, and every fiber
of $f$ is finite.
\end{lemma}

\begin{proof}
The sum of two Abelian subvarieties annihilated by $\Omega$ has the same
property, since its tangent space is the sum of their tangent spaces.
Choose one of maximal dimension; taking its sum with any other proves
both its maximality under inclusion and its uniqueness.  Complex
conjugation preserves the collection, and therefore preserves $B$.
Descent to $k$, formation of the quotient, and descent of invariant
differentials give $A$ and $\eta$.  If $D\subset A$ were a
positive-dimensional Abelian subvariety annihilated by $\eta$, the
connected inverse image of $D$ in $J$ would strictly contain $B$ and
would be annihilated by $\Omega$.  This is impossible.

Since $f^*\eta=\omega\ne0$, the map $f$ is nonconstant.  The image of
$j$ generates $J$, so $f(C)$ generates $A$.  A nonconstant morphism from
a smooth projective integral curve has finite fibers; otherwise a fiber
would contain the entire curve.  These observations also show that
$\dim A\ge1$.
\end{proof}

We call $A$ the \emph{minimal Abelian quotient of $(C,\omega)$}.  Its
dimension, rather than the genus alone or the decomposition of $J$ alone,
governs arithmetic division.  In particular, a split Jacobian does not
by itself imply that $\omega$ comes from an elliptic curve.

Our transcendence input is the following consequence of the analytic
subgroup theorem in its logarithmic form
\cite[Theorem~6.2]{lm-HuberWustholz}.  If $G/\Qbar$ is a connected
commutative algebraic group, $u\in\operatorname{Lie}(G)_{\mathbb C}$,
$\exp_G(u)\in G(\Qbar)$, and $V\subset\operatorname{Lie}(G)$ is a
$\Qbar$-linear subspace with $u\in V_{\mathbb C}$, then there is a
connected algebraic subgroup $H\subset G$, defined over $\Qbar$, such
that
\begin{equation}\label{dd-ast}
 u\in\operatorname{Lie}(H)_{\mathbb C}\subset V_{\mathbb C}.
\end{equation}
The assertion includes the case of a nonzero logarithm whose exponential
is the identity.  It is therefore applicable to periods, without an
extra non-torsion hypothesis.

\begin{lemma}[Rigidity in the kernel]\label{dd-kernel}
If $w\in\operatorname{Lie}(A)_{\mathbb C}$ satisfies
$\exp_A(w)\in A(\Qbar)$ and $\eta(w)=0$, then $w=0$.
\end{lemma}
\begin{proof}
Apply \eqref{dd-ast} with $V=\ker\eta$.  A nonzero $w$ would produce a
positive-dimensional Abelian subvariety of $A$ annihilated by $\eta$,
contrary to Lemma~\ref{dd-minimal}.
\end{proof}

Lift the path $f(P_t)$ to a continuous logarithm
$u_t\in\operatorname{Lie}(A)_{\mathbb C}$ with $u_0=0$, and put $v=u_1$.
Integration of invariant differentials gives
\begin{equation}\label{dd-log-identities}
 \exp_A(u_t)=f(P_t),\quad \exp_A(v)=f(P_1),\quad
 \eta(u_t)=tL,\quad \eta(v)=L.
\end{equation}
In the oval case $v$ is a period.  In the arc case it is a logarithm of
an algebraic point; the proofs below do not replace it by a period.

\begin{lemma}[Rational division forces vector division]\label{dd-vector}
If $q\in I\cap\mathbb Q$ and $P_q$ is algebraic, then $u_q=qv$.
Conversely, if $q\in I\cap\mathbb Q$ and $u_q=qv$, then $P_q$ is
algebraic.
\end{lemma}
\begin{proof}
For $q=m/n$, with $n>0$, the point $\exp_A(qv)$ is algebraic: its
$n$-multiple equals the algebraic point $[m]f(P_1)$, and $[n]:A\to A$
is finite.  If $P_q$ is algebraic, then $u_q-qv$ has algebraic
exponential and is annihilated by $\eta$.  Lemma~\ref{dd-kernel} gives
the required equality.  Conversely, the equality makes $f(P_q)$
algebraic.  Its finite fiber under the morphism $f$, which is defined
over $\Qbar$, consists of algebraic points.
\end{proof}

\subsection{The rational division dichotomy and the genus bounds}

\begin{theorem}[Rational division dichotomy]\label{dd-rational}
Under the preceding hypotheses, the following conditions are equivalent:
\begin{enumerate}
\item $S_{\mathbb Q}$ is infinite;
\item $S_{\mathbb Q}=I\cap\mathbb Q$;
\item $\dim A=1$;
\item there are an elliptic curve $E/\Qbar$, a nonconstant morphism
$h:C_{\Qbar}\to E$, and a nonzero invariant differential $\nu$ on $E$
such that $\omega=h^*\nu$.
\end{enumerate}
If these conditions do not hold, then
\begin{equation}\label{dd-counts}
 \#S_{\mathbb Q}\le
 \begin{cases}
 2g-2,&\text{for an oval, with }q=0\text{ included};\\
 2g-1,&\text{for an arc, with both endpoints included}.
 \end{cases}
\end{equation}
\end{theorem}

\begin{proof}
If $\dim A=1$, the nonzero functional $\eta$ has zero kernel.
Equations~\eqref{dd-log-identities} therefore imply $u_q=qv$ for every
real $q$, and Lemma~\ref{dd-vector} proves algebraicity for every
rational $q$.

Suppose instead that $\dim A>1$.  Choose a real invariant differential
$\xi$ on $A$, defined over $k$, independent of $\eta$, and put
$\theta=f^*\xi$.  The pullback on invariant differentials is injective:
it is the composite of the injective pullback from $A$ to $J$ and the
Abel--Jacobi identification with $H^0(C,\Omega^1)$.  Hence $\theta$ and
$\omega$ are independent.  Set
\[
 c=\frac{\int_\Gamma\theta}{L}\in\mathbb R,\qquad
 \delta=\theta-c\omega,\qquad
 H(t)=\int_{P_0}^{P_t}\delta.
\]
Here the path of integration is the initial part of the fixed
traversal.  The differential $\delta$ is a nonzero holomorphic
differential on $C_{\mathbb C}$, although $c$ need not be algebraic.
Consequently its zero divisor has degree $2g-2$.  The function $H$ is
real analytic, real valued, and satisfies $H(0)=H(1)=0$.  Differentiating
in the normalized integral coordinate gives
\begin{equation}\label{dd-rolle-derivative}
 H'(t)=L\,\frac{\delta}{\omega}(P_t).
\end{equation}
Thus the derivative has at most $2g-2$ zeros on a simple traversal.
The function $H$ is not identically zero, since otherwise $\delta$
would vanish on a real arc and hence on $C$.

For every $q\in S_{\mathbb Q}$, Lemma~\ref{dd-vector} implies
\[
 \int_{P_0}^{P_q}\theta=q\int_\Gamma\theta,
\]
so $H(q)=0$.  In the arc case, $N$ distinct zeros of $H$, including the
endpoints, require at least $N-1$ distinct zeros of $H'$ by Rolle's
theorem.  This gives $N\le2g-1$.  In the oval case $H$ is a real
analytic function on the parameter circle: its total increment is
zero, and the periodic differential determines its periodic analytic
extension.  Cyclic Rolle's theorem gives at least $N$ derivative zeros
from $N$ distinct zeros of $H$ on the circle.  Therefore $N\le2g-2$.
The same arguments applied to any finite subset exclude an infinite
zero set and prove \eqref{dd-counts}.

This establishes the equivalence of the first three conditions.  The
third implies the fourth by taking $E=A$ and $h=f$.  Conversely,
translate $h$ so that $h(P_0)=0$.  The identity component of the kernel
of the induced surjection $J\to E$ has codimension one and is
annihilated by $\Omega$.
Lemma~\ref{dd-minimal} then gives $\dim A\le1$, and it is nonzero.
\end{proof}

\begin{corollary}[A finite elliptic-reduction certificate]\label{dd-certificate}
On an oval, the existence of more than $2g-2$ distinct algebraic
rational-division points forces elliptic reduction of the differential
and hence algebraicity of every rational-division point.  On an arc,
the corresponding threshold is more than $2g-1$ points, with its two
endpoints counted.
\end{corollary}

For clarity, the open-arc bound permits at most $2g-3$ internal
algebraic rational-division points: its two algebraic endpoints
already account for two of the $2g-1$ points.  The closed-oval bound
counts $q=0$ once and does not count the duplicate parameter $q=1$.

\begin{remark}\label{dd-geometric-scope}
When these results are applied to polynomial lemniscates, an oval is an
oval on the smooth length cover.  A regular geometric component admits
such a traversal.  At a critical level, a geometric component may be a
graph, and a route switching analytic branches can jump between
different points of the normalization.  Such a piecewise route is not
automatically an oval or an arc in the present theorem.  The theorem
also concerns a single chosen traversal, not an unordered sum of
component lengths.
\end{remark}

\subsection{Algebraic fractions and endomorphisms}

\begin{proposition}[Rationality under scalar endomorphisms]\label{dd-scalar}
Suppose $\operatorname{End}_{\Qbar}^{0}(A)=\mathbb Q$.
If $q\in S_{\rm alg}$, then $q\in\mathbb Q$ and $u_q=qv$.
Consequently $S_{\rm alg}=S_{\mathbb Q}$, and for $\dim A>1$ the
bounds \eqref{dd-counts} hold for all algebraic fractions.
\end{proposition}

\begin{proof}
The endpoints cause no difficulty, so take $0<q<1$.  The hypothesis
implies that $A$ is absolutely simple: a nontrivial isogeny
decomposition would give a nontrivial idempotent in its rational
endomorphism algebra.  Apply \eqref{dd-ast} to $(u_q,v)$ on $A^2$ and
the algebraic hyperplane
\[
 V_q=\{(U,V):\eta(U)-q\eta(V)=0\}.
\]
The resulting connected subgroup $H$ is proper.  Both projections of
$H$ to $A$ are nonzero, since $u_q$ and $v$ are nonzero, and are thus
surjective.  Simplicity shows that $H$ has dimension $\dim A$ and each
projection is an isogeny.  As $\operatorname{End}^0(A)=\mathbb Q$,
there are nonzero integers $m,n$ such that the Lie algebra of $H$ is
the image of $W\mapsto(mW,nW)$.  Its containment in $V_q$ forces
$m-qn=0$.  Hence $q=m/n$ and $u_q=qv$.
\end{proof}

\begin{proposition}[The exact elliptic alternative]\label{dd-elliptic-alg}
If $\dim A=1$, then $S_{\rm alg}=I\cap\mathbb Q$, whether or not
$A$ has complex multiplication.
\end{proposition}
\begin{proof}
Rational fractions are already covered by Theorem~\ref{dd-rational}.
For an algebraic interior fraction with algebraic division point,
apply the preceding $A^2$ argument.  A proper connected subgroup with
both projections nonzero is the graph of a quasi-endomorphism in the
isogeny category.  Thus $q$ is the analytic scalar of a member of
$\operatorname{End}^0(A)$.  This algebra is either $\mathbb Q$ or an
imaginary quadratic field.  Its real analytic scalars are precisely
$\mathbb Q$, proving the assertion.
\end{proof}

The scalar hypothesis in Proposition~\ref{dd-scalar} should not be
dropped from its proof.  For example, on a simple Abelian variety with
real multiplication, a chosen differential may be an eigenvector for
irrational algebraic scalars.  The vector relation supplied by the
subgroup theorem can then involve a non-scalar endomorphism.  The
following result gives the general finiteness statement without
pretending that this vector relation is scalar.

\begin{lemma}[Reduction to a division hull]\label{dd-hull}
Let $Q=f(P_1)$ and
\[
 \mathcal G=\operatorname{End}_{\Qbar}(A)\cdot Q
 =\{\phi(Q):\phi\in\operatorname{End}_{\Qbar}(A)\}.
\]
This is a finitely generated subgroup of $A(\Qbar)$.  If
$q\in S_{\rm alg}$, then
\begin{equation}\label{dd-hull-eq}
 [N]f(P_q)=\phi(Q)
 \quad\text{for some }N\ge1\text{ and }
 \phi\in\operatorname{End}_{\Qbar}(A).
\end{equation}
In particular, if $Q$ is torsion, then $f(P_q)$ is torsion.
\end{lemma}

\begin{proof}
The endomorphism ring of an Abelian variety is a finitely generated
abelian group, so its image $\mathcal G$ is finitely generated.  Apply
\eqref{dd-ast} to $(u_q,v)$ and $V_q$ as above, without any simplicity
assumption, and let $H\subset A^2$ be the resulting connected Abelian
subvariety.  The identity component of the kernel of
$p_2:=\operatorname{pr}_2|_H$ lies in $A\times\{0\}$ and is
annihilated by $\eta$.  By the defining property of $A$, this identity
component is zero.  Therefore $p_2:H\to D:=p_2(H)$ is an isogeny.

In the category of Abelian varieties up to isogeny, put
$\psi=p_1p_2^{-1}\in\operatorname{Hom}^0(D,A)$.  Poincar\'e complete
reducibility supplies a rational retraction $r:A\to D$ whose
restriction to $D$ is the identity.  The quasi-endomorphism
$\Phi=\psi r\in\operatorname{End}^0(A)$ satisfies
$p_1=\Phi p_2$ on $H$ in this category.  Clear all denominators to
obtain $N\ge1$ and an actual endomorphism $\phi=N\Phi$ such that
$[N]p_1=\phi p_2$ as homomorphisms on $H$.
Since $(u_q,v)\in\operatorname{Lie}(H)_{\mathbb C}$, its exponential
$(f(P_q),Q)$ belongs to $H$.  Evaluating the last identity yields
\eqref{dd-hull-eq}.  The final assertion follows immediately.
\end{proof}

For a subgroup $\mathcal G$ of an Abelian variety, write
\[
 \operatorname{Div}(\mathcal G)
 =\{R:[N]R\in\mathcal G\text{ for some }N\ge1\}.
\]
We use the following classical curve case of the division-point
theorem: if $X$ is an integral nonelliptic projective curve in an
Abelian variety in characteristic zero and $\mathcal G$ is a finitely
generated subgroup, then
$X(\overline K)\cap\operatorname{Div}(\mathcal G)$ is finite.
The statement, including passage from a finitely generated group to
its division hull, is Theorem~1 of
\cite{dd-Oesterle1985}, following Raynaud and Faltings.  We use this
curve theorem, rather than inferring division-hull finiteness merely
from finite generation of $\mathcal G$.

\begin{theorem}[Dichotomy for all algebraic fractions]\label{dd-all-algebraic}
If $\dim A>1$, then $S_{\rm alg}$ is finite.  If $\dim A=1$, then
$S_{\rm alg}=I\cap\mathbb Q$.  Thus the following are equivalent:
infinitely many algebraic fractions have algebraic division points;
every rational fraction has an algebraic division point; and the
differential comes from an elliptic curve.
\end{theorem}

\begin{proof}
The elliptic case is Proposition~\ref{dd-elliptic-alg}.  Suppose
$\dim A>1$ and put $X=f(C)$.  This is an integral projective curve.
It is nonelliptic: a genus-one curve in an Abelian variety is a
translate of an elliptic subgroup, and if that curve contains $0$ it
is the subgroup itself.  Such an $X$ could not generate an Abelian
variety of dimension greater than one.

Lemma~\ref{dd-hull} places the images of all points counted by
$S_{\rm alg}$ in
$X\cap\operatorname{Div}(\mathcal G)$, a finite set by the preceding
division-point theorem.  The map $f$ has finite fibers, and our
traversal is injective, with its closed endpoint identified.  Hence
$S_{\rm alg}$ is finite.
\end{proof}

\begin{remark}\label{dd-effective-scope}
Theorem~\ref{dd-all-algebraic} does not provide an effective list of
the exceptional algebraic fractions, nor the numerical bounds
\eqref{dd-counts} without an additional scalar-endomorphism
hypothesis.  In the closed-oval case $Q=0$, its finiteness step is
already the Manin--Mumford theorem.  The division-hull formulation is
what also treats the open arc and non-torsion algebraic endpoints.
\end{remark}

\subsection{Optimality of the closed-oval bound in every genus}

\begin{theorem}[Sharpness]\label{dd-sharp}
For every integer $g\ge2$, there are data $(C,\omega,\Gamma,P_0)$ in
the oval setting above for which the minimal Abelian quotient has
dimension two and
\begin{equation}\label{dd-sharp-set}
 S_{\rm alg}=S_{\mathbb Q}
 =\left\{\frac{k}{2g-2}:0\le k<2g-2\right\}.
\end{equation}
Thus $2g-2$ in \eqref{dd-counts} is optimal in every genus, even if
all real algebraic fractions are allowed.
\end{theorem}

\begin{proof}
We give the algebraic base curve, the covering construction, and its
descent explicitly enough to specify an existence proof over $k$.
Start with the smooth projective model
\begin{equation}\label{dd-base-curve}
 C_0:\quad y^2=F(x),\qquad F(x)=x^5-6x^3+5x+1.
\end{equation}
The signs at $-3,-2,-1/2,1/2,2,3$ alternate as
$-,+,-,+,-,+$.  There are consequently five distinct real roots,
denoted $a_1<\cdots<a_5$.  The curve has genus two and three real
ovals.  In particular, the finite oval over $[a_1,a_2]$ carries the
positive nowhere-zero differential $\omega_0=dx/y$, with the
orientation going from $a_1$ to $a_2$ on the positive sheet.

We verify the arithmetic of its Jacobian.  Modulo $2$,
\[
 F=(x^2+x+1)(x^3+x^2+1),
\]
with both factors irreducible.  Modulo $3$, $F=x^5-x+1$ is
irreducible: it has no roots, and its remainders on division by the
three monic irreducible quadratics
$x^2+1$, $x^2+x+2$, $x^2+2x+2$ are $1$, $x+1$, $x+1$,
respectively.  Both reductions are squarefree.  Their Frobenius
cycle types give a $5$-cycle and an element of type $(2)(3)$ in the
Galois group; cubing the latter gives a transposition.  A $5$-cycle
and any transposition generate $S_5$: the conjugate transpositions
form the edge set of a connected cyclic graph on five vertices.
Hence $\operatorname{Gal}(F)=S_5$, and
\cite[Theorem~2.1]{dd-Zarhin2000} gives
\begin{equation}\label{dd-end-base}
 \operatorname{End}_{\Qbar}J(C_0)=\mathbb Z.
\end{equation}
In particular $J(C_0)$ is absolutely simple, and the minimal quotient
for $\omega_0$ is the whole two-dimensional Jacobian.

Let $d=g-1$.  The quotient of $C_0(\mathbb C)$ by conjugation is a
pair of pants $S$, namely a sphere with three boundary circles.  One
can see this directly from the three real ovals and the Euler
characteristic: the quotient has three boundary components and Euler
characteristic $-1$, which forces it to be an orientable genus-zero
surface.  The original Riemann surface is its double.

Label the boundary corresponding to the selected finite oval first.
The fundamental group of $S$ has boundary generators with product
$1$.  Send these three generators to $1,0,-1$ in $\mathbb Z/d\mathbb Z$.
This defines a connected cyclic cover $S_d\to S$ of degree $d$.
There is one boundary circle over the first boundary, $d$ over the
second, and one over the third; the first and third are degree-$d$
covers.  Since $\chi(S_d)=-d$, its $d+2$ boundary components force
its genus to be zero.  Pull back the bordered complex structure and
double along the analytic boundaries.  The resulting compact
Riemann surface $C_d$ has an antiholomorphic involution and a
holomorphic unramified map
\[
 \rho:C_d\longrightarrow C_0
\]
of degree $d$.  There is no ramification along a doubled boundary:
in a local half-disc the covering is a local biholomorphism, and
reflection extends it across the boundary.  Riemann--Hurwitz gives
$g(C_d)=d+1=g$.  The inverse image of the chosen oval is a single
real oval $\Gamma_d$ mapping onto it with degree $d$.

The construction is algebraic over $k$, not merely over $\mathbb C$.
Indeed, Riemann existence algebraizes the finite unramified cover.
The equivalence of finite \'etale cover categories under extension
of algebraically closed fields of characteristic zero descends it
from $\mathbb C$ to $\Qbar$; see
\cite[Remark~3.3(c) and Theorem~3.4]{dd-Milne2013}.
The reflection is an isomorphism from the conjugate cover to the
cover over the real base $C_0$.  Full faithfulness of the same
equivalence makes this isomorphism defined over $\Qbar$.
Its square is the identity, so it is a Galois descent datum for
$\Qbar/k$.  Effective descent for finite \'etale covers yields
$C_d$, $\rho$, and their real structure over $k$.

Set $\omega_d=\rho^*\omega_0$ and take a point above
$(a_1,0)$ as the algebraic basepoint.  The norm homomorphism
$\rho_*:J(C_d)\to J(C_0)$ is surjective.  Its connected kernel is
exactly the largest Abelian subvariety annihilated by the invariant
differential corresponding to $\omega_d$.  To check maximality,
any annihilated Abelian subvariety maps either to zero or onto the
simple $J(C_0)$.  The latter would force the nonzero differential
corresponding to $\omega_0$ to vanish identically.  It follows that
the minimal quotient is isogenous to $J(C_0)$, has dimension two,
and has rational endomorphism algebra $\mathbb Q$.

The two halves of the base oval have equal positive $\omega_0$
integrals, by the hyperelliptic involution.  The $d$ points above
each endpoint $(a_1,0)$ and $(a_2,0)$ are all on $\Gamma_d$, are
algebraic, and occur in alternating order at integral fractions
$k/(2d)$ of its total integral.  Thus there are at least $2d$
points in $S_{\mathbb Q}$.  Theorem~\ref{dd-rational} permits at
most $2g-2=2d$.  Proposition~\ref{dd-scalar} excludes additional
irrational algebraic fractions.  This proves \eqref{dd-sharp-set}.
\end{proof}

\begin{remark}\label{dd-sharp-scope}
The construction proves sharpness for holomorphic Abelian integrals
on real ovals.  It does not assert that the curves $C_d$ are length
covers of polynomial lemniscates, nor supply explicit plane
equations for every $d$.  The field of definition and the absence
of ramification are part of the proof, rather than assumptions
about a topological cover.  In particular the construction works
for both odd and even $d$.
\end{remark}

\begin{example}[The open-arc constant in genus two]\label{dd-open-sharp}
On a finite oval of the base curve \eqref{dd-base-curve}, choose an
algebraic nonbranch point $P$ and its hyperelliptic conjugate
$\iota(P)$.  The simple arc between them that passes through one
branch point has that branch point at exactly half its integral.
The endpoints and midpoint are therefore three algebraic rational
division points.  Since $g=2$ and the minimal quotient has dimension
two, the bound $2g-1=3$ excludes all others; scalar endomorphisms
also exclude irrational algebraic fractions.  This proves
sharpness of the open-arc bound in genus two.  No all-genus
sharpness assertion for open arcs is made here.
\end{example}

% DD-SECTION-DONE


\section{A complete rational-quantile classification for the beta family}
\label{bs-section}

This section classifies rational quantiles of
$\operatorname{Beta}(\alpha,1/2)$ for every positive rational $\alpha$.
The statement includes conditioning on any nondegenerate interval with
algebraic endpoints.  Its geometric form classifies division points on
monotone half-leaves of all sinusoidal spirals with positive rational
parameter.  The proof uses the analytic subgroup theorem and a cyclotomic
component of a generalized Jacobian.  In particular, it does not require
that the relevant abelian variety be simple.

\subsection{The classification and its integral form}

For $\alpha>0$ and $0\le u<v\le1$, write
\[
 F_{\alpha;u,v}(x)=
 \frac{\displaystyle\int_u^x t^{\alpha-1}(1-t)^{-1/2}\,dt}
      {\displaystyle\int_u^v t^{\alpha-1}(1-t)^{-1/2}\,dt},
 \qquad u\le x\le v.
\]
All endpoint integrals in this section are understood as convergent
improper integrals when necessary.  The function $F_{\alpha;u,v}$ is
continuous and strictly increasing from $0$ to $1$.

\begin{theorem}[Rational beta quantiles; result 76]
\label{bs-beta-classification}
Let $\alpha\in\mathbb Q_{>0}$, let $u,v\in\Qbar\cap\mathbb R$ satisfy
$0\le u<v\le1$, and let $r\in\mathbb Q\cap(0,1)$.
There is a unique $x\in(u,v)$ with $F_{\alpha;u,v}(x)=r$, and
\[
 x\in\Qbar
 \quad\Longleftrightarrow\quad
 \alpha\in\mathbb Z_{>0}\cup
 \left\{\frac12,\frac13,\frac14,\frac16\right\}.
 \tag{\ref{bs-beta-classification}.1}
\]
Consequently, outside this exceptional parameter set, every internal
rational quantile is transcendental, including every such quantile of
every algebraically truncated distribution in the family.
\end{theorem}

The equivalent integral statement is more convenient for the proof.
For coprime positive integers $a,b$, put
\[
 J_{a,b}(s)=\int_0^s\frac{t^{b-1}}{\sqrt{1-t^{2a}}}\,dt,
 \qquad 0\le s\le1.
\]
Define
\[
 \mathcal E=
 \{(1,1),(2,1),(3,1),(3,2)\}
 \ \cup\ \{(1,2k):k\ge1\}.
\]

\begin{theorem}[Monomial-weighted hyperelliptic division]
\label{bs-integral-classification}
Let $(a,b)=1$, let $0\le u<v\le1$ be algebraic, and let
$r\in\mathbb Q\cap(0,1)$.  The unique $s\in(u,v)$ satisfying
\begin{equation}
 J_{a,b}(s)-J_{a,b}(u)
 =r\bigl(J_{a,b}(v)-J_{a,b}(u)\bigr)
 \label{bs-division-equation}
\end{equation}
is algebraic if $(a,b)\in\mathcal E$, and transcendental otherwise.
\end{theorem}

The change of variables $x=s^{2a}$ gives the exact identity
\begin{equation}
 J_{a,b}(s)=\frac{1}{2a}
 B_{s^{2a}}\!\left(\frac{b}{2a},\frac12\right).
 \label{bs-beta-change}
\end{equation}
For every $\alpha\in\mathbb Q_{>0}$, the fraction
$1/(2\alpha)=a/b$ has a unique expression in coprime positive integers.
Positive real algebraicity is preserved in both directions by
$s\mapsto s^{2a}$, including at interval endpoints.  The members of
$\mathcal E$ correspond precisely to the parameter set in
Theorem~\ref{bs-beta-classification}.  Thus it suffices to prove
Theorem~\ref{bs-integral-classification}.

\subsection{Two inputs and a strict comparison}

We use the following form of the analytic subgroup theorem.  If $G$ is
a connected commutative algebraic group over $\Qbar$, $z\in
\operatorname{Lie}(G)\otimes_{\Qbar}\mathbb C$, and
$\exp_G(z)\in G(\Qbar)$, then there is a connected algebraic subgroup
$H\subseteq G$ such that
\[
 z\in\operatorname{Lie}(H)_{\mathbb C},
 \qquad
 \operatorname{Lie}(H)\subseteq\ker\ell
\]
for every algebraic linear functional $\ell$ with $\ell(z)=0$.
The case $z=0$ permits $H=0$.  This is the canonical subgroup formulation
of W\"ustholz's theorem given in
\cite[Theorem~6.2]{lm-HuberWustholz}.

We also use the standard generalized Jacobian of a smooth projective
curve $C$ with an effective modulus $\mathfrak m$.  If $S$ is the support
of $\mathfrak m$ and $P_0\notin S$, its generalized Abel map
\[
 j_{\mathfrak m}:C\setminus S\longrightarrow J_{\mathfrak m},
 \qquad P\longmapsto[P-P_0]_{\mathfrak m},
\]
is algebraic.  Pullback identifies invariant differentials on
$J_{\mathfrak m}$ with meromorphic differentials on $C$ whose poles are
bounded by $\mathfrak m$.  In characteristic zero, the connected affine
kernel of the map to the ordinary Jacobian has a torus of dimension
$|S|-1$ and a vector-group part.  These facts, including functoriality
under automorphisms preserving the modulus, are due to
Rosenlicht and Serre; see \cite{bs-rosenlicht} and
\cite[Chapter~V]{bs-serre}.  A vector group has contractible complex
points, so passing from a vector extension to its abelian quotient
induces an isomorphism on rational first homology.

\begin{lemma}[Strict comparison of weighted fractions]
\label{bs-monotonicity}
Let $0\le u<s<v\le1$ and $a,b,k$ be positive integers with $b\ne k$.
For $j=b,k$ define
\[
 R_j(s)=
 \frac{\displaystyle\int_u^s
                 t^{j-1}(1-t^{2a})^{-1/2}\,dt}
      {\displaystyle\int_u^v
                 t^{j-1}(1-t^{2a})^{-1/2}\,dt}.
\]
If $b>k$, then $R_b(s)<R_k(s)$; if $b<k$, the reverse inequality holds.
\end{lemma}

\begin{proof}
Use $t^{k-1}(1-t^{2a})^{-1/2}dt$ as a positive reference measure.
When $b>k$, the multiplier $t^{b-k}$ is strictly increasing.  Its mean
on $(u,s)$ is strictly smaller than its mean on $(s,v)$, since both
intervals have positive measure.  Multiplying out these two means gives
$R_b(s)<R_k(s)$.  Interchanging $b$ and $k$ proves the other case.
The argument remains valid when $u=0$ or $v=1$, because all the measures
in question have finite mass.
\end{proof}

\subsection{The cyclotomic component and its periods}

Assume throughout this subsection that $a\ge2$ and $(a,b)=1$.  Let
$C_a$ be the smooth projective model of
\[
 y^2=1-x^{2a},
\]
and set
\[
 \zeta=\exp(2\pi i/(2a)),\qquad
 \tau(x,y)=(\zeta x,y),\qquad
 \Omega_b=\frac{x^{b-1}\,dx}{y}.
\]
The genus is $a-1$, and its holomorphic differentials have basis
\begin{equation}
 \omega_{k-1}=\frac{x^{k-1}\,dx}{y},
 \qquad 1\le k<a,
 \qquad \tau^*\omega_{k-1}=\zeta^k\omega_{k-1}.
 \label{bs-holomorphic-basis}
\end{equation}
Together with their antiholomorphic conjugates, these characters show
that the characteristic polynomial of $\tau_*$ on
$H_1(C_a,\mathbb Q)$ is
\begin{equation}
 \frac{T^{2a}-1}{T^2-1}.
 \label{bs-homology-characteristic}
\end{equation}
In particular, every cyclotomic factor occurring here has multiplicity
one.  Put
\[
 g=\gcd(b,2a),\qquad d=\frac{2a}{g}.
\]
Coprimality of $a,b$ implies $g\in\{1,2\}$ and $d\ge3$.

\begin{lemma}[A generalized cyclotomic component]
\label{bs-generalized-component}
There exist a connected commutative algebraic group $G$ over $\Qbar$,
an algebraic map $j:C_a\setminus\{\infty_+,\infty_-\}\to G$,
and an invariant differential $\Theta_b$ on $G$ with the following
properties.
\begin{enumerate}
\item There is an exact sequence
\[
 0\longrightarrow U\longrightarrow G
   \overset{p}{\longrightarrow}A\longrightarrow0,
\]
where $U$ is a vector group and $A$ is an abelian variety.
\item $j(0,1)=0$ and $j^*\Theta_b=\Omega_b$.
\item The space $H_1(G,\mathbb Q)\cong H_1(A,\mathbb Q)$ is an
irreducible $\mathbb Q[T]$-module with characteristic polynomial
$\Phi_d(T)$, where $T$ acts by the automorphism induced by $\tau$.
\item Every $k$ with $1\le k<a$ and $\gcd(k,2a)=g$ yields an
invariant differential $\Theta_k$ pulled back from $A$ such that
$j^*\Theta_k=\omega_{k-1}$.
\item The period map
\[
 H_1(G,\mathbb Q)\longrightarrow\mathbb C,
 \qquad \gamma\longmapsto\int_\gamma\Theta_b
\]
is injective.
\end{enumerate}
\end{lemma}

\begin{proof}
The two points at infinity are exchanged by $\tau$.  Indeed, in the
coordinates $t=1/x$ and $w=y/x^a$, the two limiting values of $w$ have
opposite signs, and $\tau$ sends $w$ to $-w$.  All finite points are
regular for $\Omega_b$: at a branch point one may use $y$ as a local
parameter.  At infinity its leading order is $t^{a-b-1}dt$.
Choose an integer
$M\ge\max\{1,b-a+1\}$ and take
$\mathfrak m=M(\infty_++\infty_-)$.  This modulus is
$\tau$-invariant and bounds the poles of $\Omega_b$.

Set $J^{\mathfrak m}=J_{\mathfrak m}(C_a)$, pointed at $(0,1)$.
The latter point is fixed by $\tau$, so $\tau$ induces an automorphism
$T$ of $J^{\mathfrak m}$ of order dividing $2a$, and its Abel map is
equivariant.  The connected affine part of $J^{\mathfrak m}$ has
torus rank one.  On its torus, $T$ acts by inversion because it exchanges
the two points in the support of the modulus.  Thus this torus has
character $-1$, belonging to $\Phi_2$.

Let $e_d\in\mathbb Q[T]/(T^{2a}-1)$ be the idempotent which is $1$
on the $\Phi_d$-component and $0$ on all other cyclotomic components.
Choose a positive integer $N$ clearing its denominators, and let
\[
 f=N e_d\in\operatorname{End}(J^{\mathfrak m}),\qquad
 G=\operatorname{im}(f).
\]
Images here are taken as algebraic subgroups; $G$ is connected.
The $\Phi_2$ torus has been killed.  More explicitly, $f^2=Nf$, so
$f$ restricts to multiplication by $N$ on $G$.  Any torus in $G$ is
also a torus in $J^{\mathfrak m}$, where $f$ vanishes; it must
therefore be trivial.  Thus the connected affine part of $G$ is a
vector group.  Let $A$ be the actual
Chevalley abelian quotient of $G$.  It is isogenous to the
$\Phi_d$-factor of $\operatorname{Jac}(C_a)$.  Using the Chevalley
quotient avoids assuming that the map from $G$ to its image in the
ordinary Jacobian has connected kernel.

The invariant differential on $J^{\mathfrak m}$ pulling back to
$\Omega_b$ is a $T^*$-eigenvector of eigenvalue $\zeta^b$, which has
order $d$.  Write it as $\widetilde\Omega_b$.  If
$\iota:G\hookrightarrow J^{\mathfrak m}$ is inclusion, define
\[
 j=f\circ j_{\mathfrak m},\qquad
 \Theta_b=N^{-1}\iota^*\widetilde\Omega_b.
\]
As $e_d^*\widetilde\Omega_b=\widetilde\Omega_b$,
\[
 j^*\Theta_b
 =N^{-1}j_{\mathfrak m}^*f^*\widetilde\Omega_b
 =\Omega_b.
\]
This normalization accounts for the integer used to clear denominators.
The identical construction applies to every differential in
\eqref{bs-holomorphic-basis} of character order $d$.  Those
differentials come from the ordinary Jacobian, vanish on every affine
direction, and hence factor through $A$.

The rational first homology of $J^{\mathfrak m}$ is that of $C_a$
with the two points at infinity removed.  The extra homology class is
the toric class of character $-1$.  Consequently its $\Phi_d$-part is
the multiplicity-one $\Phi_d$-part in
\eqref{bs-homology-characteristic}.
On rational homology the image of $f$ is precisely this part, since
the scalar $N$ is invertible over $\mathbb Q$.  One may also see this
from analytic period lattices: inclusion of a connected algebraic
subgroup induces an injective map on first homology, and a surjective
homomorphism induces a surjection on rational first homology.  This
proves the asserted irreducibility and dimension $\varphi(d)$.

It remains to verify that the period map is nonzero.  Let $\delta$ be
the path from $(0,1)$ to $(1,0)$ on the positive square-root sheet,
followed by the path from $(1,0)$ to $(0,-1)$ on the negative sheet.
Then $\delta-\tau\delta$ is a closed cycle away from infinity, and
\[
 \int_{\delta-\tau\delta}\Omega_b
 =2(1-\zeta^b)J_{a,b}(1)\ne0.
\]
Its image under $j$ gives a nonzero period of $\Theta_b$.  The kernel
of the period map is $T$-stable, because
\[
 \int_{T_*\gamma}\Theta_b
   =\zeta^b\int_\gamma\Theta_b.
\]
Irreducibility over $\mathbb Q$ now implies that this kernel is zero.
\end{proof}

\begin{lemma}[Vanishing forces the abelian logarithm to vanish]
\label{bs-kernel-lemma}
In the setting of Lemma~\ref{bs-generalized-component}, suppose
$z\in\operatorname{Lie}(G)_{\mathbb C}$,
$\exp_G(z)\in G(\Qbar)$, and $\Theta_b(z)=0$.
Then $dp(z)=0$ in $\operatorname{Lie}(A)_{\mathbb C}$.
\end{lemma}

\begin{proof}
The analytic subgroup theorem gives a connected algebraic subgroup
$H\subseteq G$ with $z\in\operatorname{Lie}(H)_{\mathbb C}$ and
$\operatorname{Lie}(H)\subseteq\ker\Theta_b$.
Let $B=p(H)\subseteq A$.  It is an abelian subvariety.  The kernel of
$H\to B$ is $H\cap U$, an algebraic subgroup of a vector group and
therefore itself a vector group in characteristic zero.  It is connected
and analytically contractible.  Hence
\[
 H_1(H,\mathbb Q)\longrightarrow H_1(B,\mathbb Q)
\]
is an isomorphism.  If $\dim B>0$, inclusion of $B$ in $A$ supplies
a nonzero image in $H_1(A,\mathbb Q)$ and thus in
$H_1(G,\mathbb Q)$.  Its $\Theta_b$-period is zero, since the
restriction of this invariant differential to $H$ is zero.  This
contradicts the injectivity in Lemma~\ref{bs-generalized-component}.
Thus $B=0$, whence $H\subseteq U$ and $dp(z)=0$.
\end{proof}

\subsection{Proof of the classification}

\begin{proof}[Proof of Theorem~\ref{bs-integral-classification}]
First suppose $a\ge2$.  For $t\in[0,1]$, put
$P_t=(t,\sqrt{1-t^{2a}})$, and choose the logarithm $z_t$ of $j(P_t)$
obtained by integrating along the real path from $P_0$ to $P_t$.
The normalization in Lemma~\ref{bs-generalized-component} gives
\[
 \Theta_b(z_t)=J_{a,b}(t),\qquad
 \Theta_k(z_t)=J_{a,k}(t)
\]
for all the holomorphic indices $k$ occurring in that lemma.

Suppose that the solution $s$ of \eqref{bs-division-equation} is
algebraic, and write $r=m/h$ with positive integers $0<m<h$.
All $j(P_t)$ for $t=u,s,v$ are algebraic points.  Therefore
\[
 z=h(z_s-z_u)-m(z_v-z_u)
\]
has algebraic exponential, and \eqref{bs-division-equation} states
$\Theta_b(z)=0$.  Lemma~\ref{bs-kernel-lemma} gives $dp(z)=0$.
Every holomorphic differential $\Theta_k$ in the same cyclotomic
component consequently satisfies $\Theta_k(z)=0$.  In other words,
\begin{equation}
 \frac{J_{a,k}(s)-J_{a,k}(u)}{J_{a,k}(v)-J_{a,k}(u)}=r.
 \label{bs-forced-other-fraction}
\end{equation}
If one such index $k$ differs from $b$, this contradicts
Lemma~\ref{bs-monotonicity}.

If $b>a$, take $k=g=\gcd(b,2a)$.  For odd $b$ this is $k=1<a$;
for even $b$, the integer $a$ is odd and at least $3$, so $k=2<a$.
In either case $k<b$ and the character order is $d$.  This proves
transcendence in every case with $a\ge2$ and $b>a$, including cases
where $\Omega_b$ is of the second kind.

If $b<a$, there are exactly $\varphi(d)/2$ indices $k$ with
$1\le k<a$ and $\gcd(k,2a)=g$.  When $\varphi(d)\ge4$,
one of them differs from $b$, again giving a contradiction.  The only
remaining possibility is $\varphi(d)=2$, namely $d\in\{3,4,6\}$.
Under $b<a$ and $(a,b)=1$, these cases are exactly
\[
 (a,b)=(2,1),(3,1),(3,2).
\]

For $(2,1)$ the curve $y^2=1-x^4$ has genus one and the length
differential is a nonzero holomorphic differential.  For $(3,2)$ the
map
\[
 (x,y)\longmapsto(X,Y)=(x^2,y),\qquad Y^2=1-X^3,
\]
satisfies $dX/Y=2x\,dx/y$.  For $(3,1)$ use instead
\[
 (x,y)\longmapsto(X,Y)=(x^{-2},iy/x^3),\qquad
 dX/Y=2i\,dx/y.
\]
Both rational maps extend to nonconstant morphisms between the smooth
projective curves.  In particular the point $x=0$ in the latter map
is interpreted on the projective models.  Thus in each case
$\Omega_b$ is the pullback of a nonzero invariant differential on an
elliptic curve, up to a nonzero algebraic scalar.

On that elliptic curve the equality of scalar integrals in
\eqref{bs-division-equation} is an equality of logarithms, because
the Lie algebra is one-dimensional.  It follows that $h$ times the
image of $P_s$ is an algebraic group combination of the images of
$P_u$ and $P_v$.  Multiplication by $h$ has finite algebraic fibers,
so the image of $P_s$ is algebraic.  A nonconstant morphism of smooth
projective curves is finite, hence $P_s$, and then $s$, is algebraic.
This proves all three genus-one exceptions without requiring the
endpoints to be torsion points.

It remains to consider $a=1$.  If $b=2k$ is even, the substitution
$y=\sqrt{1-s^2}$ gives
\[
 J_{1,2k}(s)=\int_y^1(1-t^2)^{k-1}\,dt.
\]
This is a rational-coefficient polynomial in $y$.  For algebraic
endpoints, the division equation is a nonconstant polynomial equation
over $\Qbar$, so its solution $y$, and therefore $s$, is algebraic.
The same argument allows any algebraic $r\in(0,1)$.

For $b=1$, write $\theta_t=\arcsin t\in[0,\pi/2]$.
Then $J_{1,1}(t)=\theta_t$.  At an algebraic $t$,
\[
 e^{i\theta_t}=\sqrt{1-t^2}+it
\]
is algebraic.  The relation
$h(\theta_s-\theta_u)=m(\theta_v-\theta_u)$ shows that
$e^{i\theta_s}$ is an $h$-th root of an algebraic number, and hence
$s$ is algebraic.

Finally let $b=2k+1\ge3$.  Repeated use of
\begin{equation}
 J_{1,b}(t)=\frac{b-2}{b-1}J_{1,b-2}(t)
          -\frac{t^{b-2}\sqrt{1-t^2}}{b-1}
 \label{bs-odd-reduction}
\end{equation}
gives
\[
 J_{1,b}(t)=c_k\theta_t+R(t),\qquad
 c_k=4^{-k}\binom{2k}{k}>0,
\]
where $R(t)$ is algebraic at every algebraic $t\in[0,1]$.
If $s$ were algebraic, \eqref{bs-division-equation} would imply that
$D=h(\theta_s-\theta_u)-m(\theta_v-\theta_u)$ is algebraic.
Its exponential $e^{iD}$ is also algebraic.  The
Hermite--Lindemann theorem therefore forces $D=0$.
The same fraction $r$ would then be obtained from both $J_{1,b}$
and $J_{1,1}$, contrary to Lemma~\ref{bs-monotonicity}.
This proves transcendence for every remaining pair and completes the
classification.
\end{proof}

\begin{remark}[What the argument does and does not use]
\label{bs-proof-scope}
The rationality of $r$ is used when forming the integer combination of
algebraic group points.  The general theorem does not replace rational
$r$ by arbitrary algebraic $r$.  Nor does it rely on the false assertion
that every primitive cyclotomic factor of these Jacobians is simple.
The relevant irreducibility is that of a rational homology module under
one finite-order automorphism; this is a different property.  Finally,
the proof treats singular crossings of the planar spiral through its
smooth auxiliary curve and a specified monotone radial branch.
\end{remark}

\subsection{Algebraic-ratio refinements}

Although the distinction in Remark~\ref{bs-proof-scope} is substantive,
some subfamilies also admit algebraic division ratios.  Reduction of
differentials gives the following refinements.

\begin{proposition}[Algebraic-ratio refinements]
\label{bs-algebraic-ratios}
Let $0\le u<v\le1$ be algebraic and let $r\in\Qbar\cap(0,1)$.
\begin{enumerate}
\item For $a=1$ and even $b$, the solution of
\eqref{bs-division-equation} is algebraic.
\item For $a=1$ and odd $b\ge3$, that solution is transcendental.
\item Suppose $a\ge2$ and $b=2a\ell+c$ with $\ell\ge1$ and
$0<c<a$.  Then that solution is transcendental.
\end{enumerate}
The coprimality convention $(a,b)=1$ may be retained in these statements.
\end{proposition}

\begin{proof}
The polynomial substitution already proves the first assertion.
For the second, \eqref{bs-odd-reduction} implies that
\[
 (\theta_s-\theta_u)-r(\theta_v-\theta_u)
\]
is algebraic if $s$ is algebraic.  This number is an algebraic linear
form in logarithms of algebraic numbers.  Such a form is either zero
or transcendental by Baker's theorem, so it is zero; see
\cite[Proposition~15.10 and Remark~15.11]{lm-HuberWustholz}.
Lemma~\ref{bs-monotonicity}
then gives the same contradiction as before.

For the third assertion, on $y^2=1-x^{2a}$ one has
\[
 d(x^{b-2a}y)
   =(b-2a)\frac{x^{b-2a-1}\,dx}{y}
     -(b-a)\frac{x^{b-1}\,dx}{y}.
\]
Iterating this identity gives
\begin{equation}
 J_{a,b}(t)=C J_{a,c}(t)+S(t),
 \qquad C\in\mathbb Q_{>0},
 \label{bs-general-reduction}
\end{equation}
where $S(t)$ is algebraic at algebraic $t$.  The differential defining
$J_{a,c}$ is holomorphic because $c<a$.  If the proposed division point
were algebraic, the corresponding algebraic linear combination of its
holomorphic integrals would be algebraic, hence zero.  The latter fact
is the holomorphic case of the analytic subgroup theorem; it also
follows from \cite[Proposition~13.11]{lm-HuberWustholz}, by placing
the finitely many paths on a product of Jacobians and absorbing their
algebraic coefficients into the invariant differential.
Thus the normalized partial integrals with weights $b$ and $c$ would
both equal $r$, contradicting Lemma~\ref{bs-monotonicity}.
\end{proof}

\subsection{Sinusoidal spirals and Erd\H{o}s lemniscates}

Let $q=a/b>0$ be reduced.  On a chosen radial half-leaf of the
sinusoidal spiral
\[
 r^q=2\cos(q\theta),
\]
the radius increases from $0$ to $2^{1/q}$.  Direct differentiation
of the polar equation gives
\[
 ds_{\mathrm{arc}}=
 \frac{dr}{\sqrt{1-r^{2q}/4}}.
\]
On putting $r=2^{b/a}t^b$, its length from the origin to parameter $s$
is
\begin{equation}
 b\,2^{b/a}J_{a,b}(s).
 \label{bs-spiral-length}
\end{equation}
An endpoint on this branch is algebraic if and only if its radius,
equivalently its parameter $s$, is algebraic.  For the nontrivial
direction, if $r$ is algebraic then $e^{iq\theta}$ satisfies a
quadratic equation over $\Qbar$, and
$(e^{i\theta})^a=(e^{iq\theta})^b$ makes $e^{i\theta}$ algebraic.

\begin{corollary}[All positive rational spiral parameters]
\label{bs-spiral-classification}
For a positive rational sinusoidal spiral, a nontrivial rational
length division of a monotone half-leaf, or of any subarc of that
half-leaf with algebraic endpoints, has algebraic coordinates exactly
when
\[
 q\in\left\{1,\frac32,2,3\right\}
   \cup\left\{\frac1{2k}:k\ge1\right\}.
\]
For every other positive rational $q$, the radius of every such
division point is transcendental, and the point is not algebraic.
\end{corollary}

\begin{corollary}[Division of the integer-petal lemniscates; result 73]
\label{bs-erdos-divisions}
For $n\ge1$, consider
\[
 C_n=\{z\in\mathbb C:|z^n-1|=1\}.
\]
Every internal rational length division point on a monotone half-leaf
is algebraic for $n=1,2,3$, and has transcendental radius for every
$n\ge4$.  The same assertions hold for subarcs with algebraic
endpoints contained in one monotone half-leaf.

For $n\ge4$, both Cartesian coordinates of every such nonalgebraic
division point are separately transcendental.  In a cyclic traversal
of all petals starting at the origin, the $j$-th position of an
$m$-fold equal-length division, $0\le j<m$, is algebraic exactly when
\[
 m\mid 2nj.
\]
These exceptional positions are the origin and petal tips.
There are exactly $\gcd(m,2n)$ algebraic indexed positions among the
$m$ positions; all $m$ positions are algebraic if and only if $m\mid2n$.
\end{corollary}

\begin{proof}
Here $q=n$, and \eqref{bs-spiral-length} has $(a,b)=(n,1)$.
This proves the half-leaf statement.  The real equation
\[
 (X^2+Y^2)^n-2\operatorname{Re}(X+iY)^n=0
\]
is monic of degree $2n$ as a polynomial in either coordinate, with
coefficients polynomial in the other.  If one coordinate of a point
were algebraic, the other and its radius would also be algebraic.

There are $2n$ equal-length half-leaves in one traversal.  Its division
position $j/m$ lies at a half-leaf endpoint precisely when
$2nj/m$ is an integer.  Otherwise it corresponds, after reversing the
orientation if necessary, to an internal rational fraction of one
half-leaf, to which the first statement applies.
The congruence $2nj\equiv0\pmod m$ has exactly $\gcd(m,2n)$ solutions
modulo $m$, proving the two counting assertions.
\end{proof}

The final counting statement concerns arclength positions: the same
geometric origin is visited between consecutive petals.  It is not a
claim that all those indexed positions are distinct planar points.
The midpoint of a complete leaf is its algebraic tip, which is why
``half-leaf'' cannot be omitted from the nontrivial-division assertion.
Tamiozzo established the cases $n\le3$ and asked about $n>3$ in
\cite[Remark~5.7]{lm-Tamiozzo}.  Corollary~\ref{bs-erdos-divisions}
answers that specific question through the rational-quantile
classification above.  No assertion of historical priority is made.

\subsection{Polynomial images of the circle}

An explicit contrast is supplied by
\[
 z(t)=(1+e^{it})^m,\qquad -\pi\le t\le\pi,
 \qquad m\in\mathbb Z_{>0}.
\]
On either monotone radial half of this parametrization,
$r=2^m\cos^m(t/2)$ and $\theta=mt/2$; thus the corresponding
spiral parameter is $q=1/m$.  Its total length is
\[
 L_m=m\,2^{m+1}J_{1,m}(1).
\]
For even $m$ this is rational.  For $m=2k+1$ it is
\[
 L_{2k+1}=2(2k+1)\binom{2k}{k}\pi.
\]
The finite self-intersections do not affect these lengths: each
monotone radial half is injective, and any intersections between the
two halves are isolated.

\begin{example}[A parity contrast]
\label{bs-circle-powers}
For even $m$, every internal algebraic length ratio on a radial
half-arc with algebraic endpoints produces an algebraic point.  For
odd $m\ge3$, every such ratio produces a point with transcendental
radius, by Proposition~\ref{bs-algebraic-ratios}.
The first few total lengths are
\[
 L_1=2\pi,\qquad L_2=16,\qquad L_3=12\pi,
 \qquad L_4=\frac{256}{3},\qquad L_5=60\pi.
\]
In particular, being an algebraic multiple of $\pi$ does not ensure
algebraicity of equal-arclength division points.
\end{example}

\begin{example}[The cubic image]
\label{bs-cubic-image}
For $z=(1+e^{it})^3$, the radius is $8s^3$.  The midpoint of its
monotone radial half-arc is specified by the unique $s\in(0,1)$ with
\[
 \arcsin s-s\sqrt{1-s^2}=\frac\pi4.
\]
This $s$ is transcendental, as is $8s^3$.  Equivalently,
$x=s^2$ is the median of $\operatorname{Beta}(3/2,1/2)$.
The planar image has equation
\[
 (X^2+Y^2-2X)^3=27(X^2+Y^2)^2.
\]
This equation is monic of degree six in either Cartesian coordinate.
Consequently both coordinates of the indicated midpoint are
transcendental: algebraicity of one would imply algebraicity of the
other and of the radius.
\end{example}

These conclusions concern exact arithmetic values and exact division
points.  They neither obstruct numerical approximation nor imply
algebraic independence of the different quantiles.


\section{Algebraic levels and torsion packets for generalized lemniscatic integrals}
\label{tr-sec}

For an integer $n\geq 2$, put
\begin{equation}\label{tr-def}
 I_n(s)=\int_0^s\frac{dt}{\sqrt{1-t^{2n}}},\qquad
 F_n(s)=\frac{I_n(s)}{I_n(1)}\qquad(0\leq s\leq 1).
\end{equation}
The square root is positive on $(0,1)$.  The endpoint singularity is
integrable, and $F_n$ is a continuous strictly increasing bijection of
$[0,1]$ onto itself.  A rational division problem asks whether
$F_n^{-1}(q)$ is algebraic for $q\in\mathbb Q$.  In this section we allow
\emph{every algebraic real level} $q$.  The new input relative to rational
division is a torsion consequence of the analytic subgroup theorem,
followed by explicit, previously established classifications of torsion
points on curves.  The torsion classifications themselves are not new
results of this paper.

\begin{theorem}\label{tr-main}
Let $n=4$, or let $n\geq 5$ be an odd prime.  Then
\[
 s\in\Qbar\cap(0,1)\quad\Longrightarrow\quad F_n(s)\notin\Qbar.
\]
Equivalently, for every $q\in\Qbar\cap(0,1)$, the unique real solution
$s_{n,q}\in(0,1)$ of
\begin{equation}\label{tr-quantile}
 \int_0^{s_{n,q}}\frac{dt}{\sqrt{1-t^{2n}}}
 =q\int_0^1\frac{dt}{\sqrt{1-t^{2n}}}
\end{equation}
is transcendental.
\end{theorem}

The equivalence in the theorem uses only strict monotonicity.  Notice
that neither assertion follows from the transcendence of the two
integrals separately: a quotient of transcendental numbers can be
algebraic.  The proof will exclude that possibility through the
arithmetic of the associated Abel--Jacobi point.

\subsection{A torsion criterion for an algebraic ratio}
\label{tr-subsec-criterion}

Let $A/\Qbar$ be an abelian variety and let $\eta$ be a nonzero invariant
algebraic differential on $A$.  We also write $\eta$ for its linear
functional on $\operatorname{Lie}(A)$.  Call $(A,\eta)$
\emph{minimal for the differential} if no positive-dimensional abelian
subvariety $B\subset A$ satisfies $\eta|_B=0$.

This condition can always be obtained by taking a quotient.  Indeed,
the sum of two abelian subvarieties annihilated by $\eta$ is again
annihilated by $\eta$.  Finite-dimensionality therefore gives a largest
such connected abelian subvariety $B$.  The differential descends to
$A/B$, and the descended pair is minimal.  All these objects and maps
are defined over $\Qbar$.

We use the following form of the analytic subgroup theorem.  If $G$
is a connected commutative algebraic group over $\Qbar$, $z\in
\operatorname{Lie}(G)_{\mathbb C}$ has algebraic exponential, and an
algebraic linear subspace $V\subset\operatorname{Lie}(G)$ contains $z$
after scalar extension, then there is a connected algebraic subgroup
$H\subset G$ such that
\[
 z\in\operatorname{Lie}(H)_{\mathbb C}
 \subset V_{\mathbb C}.
\]
This follows directly from the annihilator formulation in
\cite[Theorem~6.2]{lm-HuberWustholz}.

\begin{lemma}[Algebraic ratio to a torsion logarithm]\label{tr-torsion-lemma}
Suppose $(A,\eta)$ is minimal for the differential.  Let
$u,v\in\operatorname{Lie}(A)_{\mathbb C}$ satisfy
\[
 \exp_A(u)\in A(\Qbar),\qquad
 \exp_A(v)\in A(\Qbar)_{\mathrm{tors}}.
\]
If $q\in\Qbar^{\times}$ and $\eta(u)=q\eta(v)$, then
$\exp_A(u)$ is a torsion point.
\end{lemma}

\begin{proof}
Apply the analytic subgroup theorem to $G=A\times A$, the logarithm
$(u,v)$, and the algebraic hyperplane
\[
 V_q=\{(U,V):\eta(U)-q\eta(V)=0\}.
\]
It gives a connected abelian subvariety $H\subset A^2$ containing
$(u,v)$ in its complex Lie algebra and having Lie algebra contained
in $V_q$.  Let $K$ be the identity component of the kernel of the
second projection $\operatorname{pr}_2:H\to A$.  It has the form
$B\times\{0\}$ for an abelian subvariety $B\subset A$.  The hyperplane
condition implies $\eta|_B=0$.  Minimality forces $B=0$, so
$\operatorname{pr}_2|_H$ has finite kernel.

The point
$h=(\exp_A(u),\exp_A(v))$ belongs to $H(\mathbb C)$.  Choose an integer
$m\geq1$ killing its second coordinate.  Then $[m]h$ lies in the finite
kernel of $\operatorname{pr}_2|_H$, hence is torsion.  Consequently $h$
is torsion, and so is its first coordinate.
\end{proof}

\begin{remark}\label{tr-lemma-scope}
The lemma does not require $A$ to be simple, and it does not require
$v$ to be a period: a logarithm of any torsion point suffices.  These
two features are useful here.  In particular, the genus-two Jacobian
occurring for $n=4$ must not be treated as a simple abelian variety.
Minimality of the differential is the property actually needed.
\end{remark}

\subsection{The curve carrying the differential and its minimal quotient}
\label{tr-subsec-curve}

Let $C_n$ be the smooth projective model of
\begin{equation}\label{tr-curve}
 v^2=u(1-u^n).
\end{equation}
The polynomial on the right is squarefree, and
\begin{equation}\label{tr-genus}
 g(C_n)=\left\lfloor\frac n2\right\rfloor.
\end{equation}
The differential
\begin{equation}\label{tr-omega}
 \omega_n=\frac{du}{2v}
\end{equation}
is a nonzero holomorphic differential on $C_n$.  For $0\leq s\leq1$,
consider the real path
\begin{equation}\label{tr-path}
 P(s)=\bigl(s^2,s\sqrt{1-s^{2n}}\bigr).
\end{equation}
It runs from $W_0=(0,0)$ to $W_1=(1,0)$, and substitution gives
\begin{equation}\label{tr-pullback}
 \int_{W_0}^{P(s)}\omega_n=I_n(s).
\end{equation}
The integrals are taken along this specified path.  Although the
displayed affine formula uses branch-point endpoints, the path and
the differential are regular in local parameters on the smooth curve.

Both $W_0$ and $W_1$ are Weierstrass points.  In particular,
\begin{equation}\label{tr-endpoint-torsion}
 \operatorname{div}\left(\frac{u-1}{u}\right)=2W_1-2W_0,
 \qquad [W_1-W_0]\in J(C_n)[2].
\end{equation}
Let $j_n:C_n\to J_n=J(C_n)$ be the Abel--Jacobi map based at $W_0$,
and let $\Omega_n$ be the invariant differential with
$j_n^*\Omega_n=\omega_n$.  Denote its minimal quotient by $A_n$.

\begin{proposition}[Cyclotomic lower bound]\label{tr-cyclotomic}
For every $n\geq2$,
\[
 \dim A_n\geq\frac{\varphi(2n)}2.
\]
In particular, $A_n=J_n$ when $n=4$ or when $n$ is an odd prime.
Here the equality means that the largest connected abelian subvariety
annihilated by $\Omega_n$ is zero.
\end{proposition}

\begin{proof}
Choose $\zeta_{2n}=e^{\pi i/n}$ and set $\zeta_n=\zeta_{2n}^2$.
The automorphism
\[
 \rho(u,v)=(\zeta_nu,\zeta_{2n}v)
\]
of $C_n$ fixes $W_0$ and satisfies
\begin{equation}\label{tr-eigenform}
 \rho^*\omega_n=\zeta_{2n}\omega_n.
\end{equation}
It therefore induces an automorphism of $J_n$ with the same
eigenvalue on $\Omega_n$.  If an abelian subvariety is annihilated by
$\Omega_n$, its image under $\rho$ is also annihilated.  The largest
such subvariety is thus $\rho$-stable, and $\rho$ descends to $A_n$.

On $A_n$, the descended holomorphic differential is still an
eigenvector with eigenvalue the primitive $2n$-th root $\zeta_{2n}$.
The action on rational first cohomology has dimension $2\dim A_n$
and rational characteristic polynomial.  Its complexification contains
the holomorphic cohomology and hence the eigenvalue $\zeta_{2n}$.
Its characteristic polynomial must therefore be divisible by
$\Phi_{2n}$, of degree $\varphi(2n)$.  This proves the lower bound.

For $n=4$, the bound is $\varphi(8)/2=2=g(C_4)$.  For an odd prime
$p$, it is $\varphi(2p)/2=(p-1)/2=g(C_p)$.  Since $A_n$ is a quotient
of $J_n$, equality of dimensions forces the connected kernel to be
zero in these cases.
\end{proof}

This argument deliberately avoids a simplicity assertion for $J(C_4)$.
The simple cyclotomic factors for odd primes are also described in
\cite[Proposition~4.8]{lm-Tamiozzo}, but their simplicity is not needed
in the proof just given.

\begin{corollary}\label{tr-reduction-to-packet}
Suppose $n=4$ or $n\geq5$ is an odd prime.  If $s\in\Qbar\cap(0,1)$
and $F_n(s)\in\Qbar$, then $[P(s)-W_0]$ is torsion in $J_n$.
\end{corollary}

\begin{proof}
Let $u_s$ and $v_1$ be the Abel logarithms obtained by integrating from
$W_0$ to $P(s)$ and $W_1$, respectively, along the chosen path.
Their exponentials are $[P(s)-W_0]$ and $[W_1-W_0]$.  The former is
algebraic and the latter is torsion by \eqref{tr-endpoint-torsion}.
Moreover,
\[
 \Omega_n(u_s)=I_n(s)=F_n(s)I_n(1)=F_n(s)\Omega_n(v_1).
\]
The coefficient $F_n(s)$ is positive.  Proposition~\ref{tr-cyclotomic}
makes $(J_n,\Omega_n)$ minimal, so Lemma~\ref{tr-torsion-lemma} applies.
\end{proof}

\subsection{Excluding the real interval from the torsion packets}
\label{tr-subsec-packets}

\begin{proof}[Proof of Theorem~\ref{tr-main} for $n=4$]
Write $\zeta_8=e^{\pi i/4}$ and choose $c\in\Qbar$ with
$c^2=\zeta_8$.  The transformation
\begin{equation}\label{tr-bolza-transform}
 X=\zeta_8u,\qquad Y=cv
\end{equation}
identifies $C_4:v^2=u-u^5$ with the smooth projective curve
\begin{equation}\label{tr-bolza}
 B:Y^2=X^5+X.
\end{equation}
Indeed $X^5+X=\zeta_8(u-u^5)=Y^2$.  Let $\infty$ denote the unique
point at infinity on this odd-degree model.

Boxall and Grant classify the entire torsion packet of $B$ based at
$\infty$ in \cite[Theorem~(0.3)]{tr-boxall-grant}.  It consists of
$\infty$, the five finite Weierstrass points, and sixteen further
points.  The possible finite $X$-coordinates are exactly the roots of
\begin{equation}\label{tr-bolza-packet}
 X^5+X=0,\qquad X^4+4X^2+1=0,\qquad X^4-4X^2+1=0.
\end{equation}
The sixteen non-Weierstrass points have order six in this embedding.
The reference uses $\infty$ as base point, whereas our base point maps
to $(0,0)$.  This change does not affect the torsion condition: the
function $X$ has divisor $2(0,0)-2\infty$, so
$[(0,0)-\infty]$ is a $2$-torsion point.

At $P(s)$ with $0<s<1$, we have $u=s^2\in(0,1)$ and
$X=\zeta_8u$.  The first equation in \eqref{tr-bolza-packet} is
impossible because $X\ne0$ and $X^4+1=1-u^4>0$.  The other two
equations become
\[
 -u^4+4iu^2+1=0,\qquad -u^4-4iu^2+1=0,
\]
both impossible because their imaginary parts are nonzero.  Thus no
point on the open path belongs to this torsion packet.  If $s$ and
$F_4(s)$ were both algebraic, Corollary~\ref{tr-reduction-to-packet}
would give precisely such a torsion point.  This proves the assertion.
\end{proof}

\begin{proof}[Proof of Theorem~\ref{tr-main} for odd primes $p\geq5$]
The transformation
\begin{equation}\label{tr-prime-transform}
 X=-\frac1u,\qquad Y=\frac{iv}{u^{(p+1)/2}}
\end{equation}
identifies the smooth projective model of $C_p$ with
\begin{equation}\label{tr-prime-curve}
 D_p:Y^2=X^p+1.
\end{equation}
Indeed both $Y^2$ and $X^p+1$ equal $(u^p-1)/u^p$.  The inverse is
rational, and the birational map extends to an isomorphism of smooth
projective curves.  The base point $W_0$ maps to the unique point
$\infty$ of $D_p$: $u$ has a double zero at $W_0$, so $X$ has a double
pole there.  Hence the base point is exactly the one used in the
following torsion classification.

Arul's classification \cite[Theorem~6.2.21]{tr-arul} gives the whole
torsion packet of $D_p$ based at $\infty$.  For $p\geq7$, its points
are precisely $\infty$ and the points with $X=0$ or $Y=0$.  For $p=5$
there are, in addition, the ten points
\begin{equation}\label{tr-p5-exceptions}
 (X,Y)=\bigl(4^{1/5}\zeta_5^j,\ \pm\sqrt5\bigr),
 \qquad j=0,1,2,3,4.
\end{equation}
The latter case is due to Coleman and is included with attribution
in Arul's theorem.  These are statements about the complete geometric
torsion packet, not just points of small order or points over a fixed
number field.

For $P(s)$ with $0<s<1$, formula \eqref{tr-prime-transform} reads
\begin{equation}\label{tr-prime-real-image}
 X=-s^{-2}<-1,\qquad
 Y=i\frac{\sqrt{1-s^{2p}}}{s^p}\in i\mathbb R_{>0}.
\end{equation}
It is therefore neither at infinity nor a point with $X=0$ or $Y=0$.
For $p=5$, its nonzero purely imaginary $Y$-coordinate also excludes
every point in \eqref{tr-p5-exceptions}.  Consequently no point on
the open path is torsion relative to $W_0$.  Applying
Corollary~\ref{tr-reduction-to-packet} proves the theorem.
\end{proof}

\subsection{Explicit levels and geometric division}
\label{tr-subsec-consequences}

For example, each of the two equations
\begin{align}
 \int_0^a\frac{dt}{\sqrt{1-t^8}}
 &=\frac1{\sqrt2}\int_0^1\frac{dt}{\sqrt{1-t^8}},
 \label{tr-example-eight}\\
 \int_0^b\frac{dt}{\sqrt{1-t^{10}}}
 &=\frac{\sqrt5-1}{2}\int_0^1\frac{dt}{\sqrt{1-t^{10}}}
 \label{tr-example-ten}
\end{align}
has a unique solution in $(0,1)$, and both solutions are
transcendental.  These examples use algebraic irrational levels, so
they are not instances of a rational-division conclusion alone.

There is also a direct interpretation on the Erd\H{o}s lemniscate
\begin{equation}\label{tr-lemniscate}
 \Lambda_n=\{z\in\mathbb C:|z^n-1|=1\}.
\end{equation}
On the half-leaf in the first quadrant joining the origin to
$2^{1/n}$, polar coordinates satisfy
$r^n=2\cos(n\theta)$.  If $r=2^{1/n}s$, the arclength from the origin
is $2^{1/n}I_n(s)$; hence its fraction of the half-leaf's total
arclength is exactly $F_n(s)$.

\begin{corollary}\label{tr-geometric-cor}
Let $n=4$ or let $n\geq5$ be an odd prime.  For every algebraic
$q\in(0,1)$, the point at arclength fraction $q$ on this half-leaf has
transcendental radius, transcendental real coordinate, and
transcendental imaginary coordinate.
\end{corollary}

\begin{proof}
Its radius is $2^{1/n}s_{n,q}$ and is transcendental by
Theorem~\ref{tr-main}.  Write the point as $x+iy$.  If either $x$ or
$y$ were algebraic, the equation
\[
 \bigl((x+iy)^n-1\bigr)\bigl((x-iy)^n-1\bigr)=1
\]
would be a nonconstant polynomial equation with algebraic
coefficients for the other coordinate.  Its leading coefficient in
that coordinate is $1$.  Both coordinates would then be algebraic,
and so would $\sqrt{x^2+y^2}$, a contradiction.
\end{proof}

\begin{remark}[Scope of the classification]\label{tr-scope}
The theorem covers $n=4$ and every odd prime at least five; it does
not assert the same conclusion for all composite $n$.  For other
$n$, a torsion conclusion in the minimal differential quotient need
not imply torsion in the entire Jacobian, and a torsion-packet theorem
for the entire curve cannot simply be applied to that quotient.
The proof above isolates and verifies this issue before invoking the
packet classifications.  Its conclusion is transcendence of each
specified quotient or inverse value, not algebraic independence of
different values.  No claim of historical priority for these
corollaries is made here.
\end{remark}


\part{Explicit independent Abelian integrals}

\section{Independent hyperelliptic integrals}
\label{ii-sec-integrals}

The passage from an algebraic curve to its Jacobian permits a stronger
conclusion than the transcendence of a single positive integral.  In this
section we give explicit families of incomplete integrals whose coordinates
are linearly independent over $\Qbar$.  There are two distinct constructions.
The first uses positivity of a Vandermonde determinant and works in every
genus, for at most $g$ ordered intervals on a genus-$g$ curve.  The second
uses independent quadratic characters and a local formal-group argument; it
produces infinitely many independent integrals on one fixed genus-two
curve.  The transcendence input is the analytic subgroup theorem.  The
endomorphism-ring input is an existing theorem of Zarhin.  These inputs,
rather than a new transcendence theorem for arbitrary periods, account for
the conclusions below.

Throughout this section, a choice of logarithm of a nonzero algebraic
number means any complex logarithm.  Set
\begin{equation}\label{ii-eq-log-space}
 \mathcal L_{\rm alg}
 =\Qbar+
   \operatorname{span}_{\Qbar}
       \{\lambda\in\mathbb C:e^\lambda\in\Qbar^{\times}\}.
\end{equation}
In particular, $1$ and $\pi$ belong to this space.  Our separation statements
concern a vector space of logarithms, not the field generated by logarithms.

\subsection{The analytic subgroup input and a separation lemma}

We use the following precise form of the analytic subgroup theorem
\cite[Theorem~6.2]{lm-HuberWustholz}.  If $G/\Qbar$ is a connected
commutative algebraic group and $u\in\operatorname{Lie}(G)_{\mathbb C}$
satisfies $\exp_G(u)\in G(\Qbar)$, there is a connected algebraic subgroup
$H\subseteq G$ such that
\[
 u\in\operatorname{Lie}(H)_{\mathbb C},\qquad
 \operatorname{Ann}_{\Qbar}(u)
   =\operatorname{Ann}_{\Qbar}(\operatorname{Lie}(H)).
\]
Here an annihilator is taken in the algebraic cotangent space of $G$.
Consequently, every algebraic linear form that vanishes at $u$ vanishes
on $\operatorname{Lie}(H)$.  The inclusion of the particular logarithm $u$
in this Lie algebra is important: one must not replace it merely by a
statement about the algebraic point $\exp_G(u)$.

\begin{lemma}[Logarithms on a power of a simple abelian variety]
\label{ii-lem-logarithms}
Let $A/\Qbar$ have dimension $g$ and satisfy
$\operatorname{End}_{\Qbar}(A)=\mathbb Z$.  Fix an algebraic basis of
$\operatorname{Lie}(A)$, and write $u_j=(u_{1j},\ldots,u_{gj})$ for
$1\le j\le r$.  Suppose that
\[
 \exp_A(u_j)\in A(\Qbar)
 \quad\hbox{and}\quad
 \sum_{j=1}^r m_j u_j=0,\ m_j\in\mathbb Z
       \ \Longrightarrow\ m_1=\cdots=m_r=0.
\]
Then all $gr$ coordinates $u_{kj}$ are linearly independent over $\Qbar$.
More precisely,
\begin{equation}\label{ii-eq-separated}
 \sum_{j=1}^r\sum_{k=1}^g c_{kj}u_{kj}\in\mathcal L_{\rm alg},
 \qquad c_{kj}\in\Qbar,
 \quad\Longrightarrow\quad c_{kj}=0\ \hbox{for all }k,j.
\end{equation}
\end{lemma}

\begin{proof}
First, $A$ is absolutely simple.  Indeed, a proper nonzero abelian
subvariety would, by Poincare reducibility, yield a nontrivial idempotent in
$\operatorname{End}^0(A)=\mathbb Q$.  We also need the following standard
description of proper abelian subvarieties of $A^r$.  If $B\subsetneq A^r$
is such a subvariety, then a nonzero homomorphism $A^r\to A$ vanishes on
$B$.  To see this, pass to the isogeny category, where the quotient
$A^r/B$ is a nonzero power of $A$, project to one factor, and clear
denominators.  Since $\operatorname{End}(A)=\mathbb Z$, the homomorphism
has the form
\[
 (a_1,\ldots,a_r)\longmapsto\sum_j[m_j]a_j
\]
for a nonzero integer row $(m_1,\ldots,m_r)$.
It follows from the hypothesis on the $u_j$ that no proper abelian
subvariety of $A^r$ has Lie algebra containing $u=(u_1,\ldots,u_r)$.

Suppose now that the left side of \eqref{ii-eq-separated} equals
$c+\sum_{h=1}^s b_h\lambda_h$, where $c,b_h\in\Qbar$ and
$e^{\lambda_h}=\alpha_h\in\Qbar^{\times}$.  Consider
\[
 G=A^r\times\mathbb G_a\times\mathbb G_m^s,
 \qquad v=(u,1,\lambda_1,\ldots,\lambda_s).
\]
Its exponential is algebraic.  The assumed equality supplies an
algebraic linear form $\ell$ with $\ell(v)=0$.  Apply the stated analytic
subgroup theorem and let $H$ be the resulting connected subgroup.

At the level of Lie algebras, a connected algebraic subgroup of a product
of an abelian variety and a commutative affine algebraic group splits
between these two factors.  One way to see this is to apply the subgroup
version of Goursat's lemma: any relation between the factors would give a
common connected quotient, simultaneously proper and affine, hence
trivial.  Finite kernels do not affect Lie algebras.  Thus, if $B$ is the
abelian projection of $H$, then
$\operatorname{Lie}(B)\times\{0\}\subseteq\operatorname{Lie}(H)$.
As $u\in\operatorname{Lie}(B)_{\mathbb C}$, the preceding paragraph
forces $B=A^r$.  The form $\ell$ must therefore vanish on the entire
$A^r$ direction.  In the chosen algebraic basis, this says exactly that
all the $c_{kj}$ vanish.
\end{proof}

\begin{remark}\label{ii-rem-limits}
The hypothesis is independence of the logarithm vectors over
$\mathbb Z$, not algebraic independence of their coordinates.  The
conclusion is also linear independence, not algebraic independence.
For example, a determinant identity involving products of integrals is
not decided by this lemma.  Separation from
$\mathcal L_{\rm alg}$ does imply that the exponential of any nonzero
algebraic linear combination of the coordinates is transcendental.
If $U$ is such a combination, $e^U\in\Qbar$ would place
$U$ in \eqref{ii-eq-log-space}.  Applying this observation to $iU$ also
shows that $\sin U$ and $\cos U$ are transcendental: an algebraic value
of either would make $e^{iU}$ a root of a quadratic polynomial over
$\Qbar$.
\end{remark}

\subsection{A positive determinant in every genus}

For $g\ge2$, put
\begin{equation}\label{ii-eq-genus-family}
 f_g(x)=x^{2g+1}-x+1,\qquad
 C_g:y^2=f_g(x),\qquad J_g=\operatorname{Jac}(C_g),
\end{equation}
where $C_g$ denotes the smooth projective model.  These curves have genus
$g$.  On $[0,1]$ one has $f_g(x)=1-x+x^{2g+1}>0$.  We always take the
positive square root there.

\begin{proposition}[The endomorphism input]
\label{ii-prop-endomorphisms}
For every $g\ge2$,
\[
 \operatorname{End}_{\Qbar}(J_g)=\mathbb Z.
\]
\end{proposition}
\begin{proof}
Zarhin's theorem states that in characteristic zero the Jacobian of
$Y^2=h(X)$ has endomorphism ring $\mathbb Z$ when $\deg h\ge5$ and
the Galois group of $h$ is the full symmetric or alternating group.
His Example~2.2(1) explicitly applies this to
$h(X)=X^n-X-1$ for every $n\ge5$
\cite[Theorem~2.1 and Example~2.2(1)]{dd-Zarhin2000}.
For odd $n=2g+1$, the change of variables
\[
 X=-x,\qquad Y=iy
\]
identifies $C_g$ over $\Qbar$ with that curve.  Endomorphism rings are
unchanged by this isomorphism.  Thus no unverified assertion about a
generic specialization is needed for \eqref{ii-eq-genus-family}.
\end{proof}

\begin{lemma}[An interval moment determinant]
\label{ii-lem-vandermonde}
Let $w$ be positive and continuous on $[0,1]$, and let
\[
 0\le a_1<b_1\le a_2<b_2\le\cdots\le a_r<b_r\le1.
\]
Then
\begin{equation}\label{ii-eq-vandermonde}
 \det\left(\int_{a_j}^{b_j} t^{k-1}w(t)\,dt\right)_{1\le k,j\le r}
 =\int_{\prod_j[a_j,b_j]}
   \prod_{i<j}(t_j-t_i)\prod_{j=1}^r w(t_j)\,
   dt_1\cdots dt_r>0.
\end{equation}
\end{lemma}
\begin{proof}
Multilinearity of the determinant and Fubini's theorem give the equality.
The determinant of $(t_j^{k-1})_{k,j}$ is the displayed Vandermonde
product.  Every factor is nonnegative on the integration box, and all
factors are strictly positive in its interior.  The box has positive
measure and $w$ is positive, proving the strict inequality.  Common
endpoints between adjacent intervals have measure zero and cause no
problem.
\end{proof}

\begin{theorem}[Ordered interval integrals; result 72]
\label{ii-thm-intervals}
Let $g\ge2$, let $1\le r\le g$, and let the endpoints in
Lemma~\ref{ii-lem-vandermonde} be real algebraic numbers.  Define
\begin{equation}\label{ii-eq-block-integrals}
 M_{k,j}=\int_{a_j}^{b_j}
             \frac{t^k\,dt}{\sqrt{t^{2g+1}-t+1}},
 \qquad 0\le k\le g-1,\quad1\le j\le r.
\end{equation}
All $gr$ numbers $M_{k,j}$ are linearly independent over $\Qbar$,
even modulo $\mathcal L_{\rm alg}$ in the sense of
\eqref{ii-eq-separated}.  In particular, the family consisting of
$1$, $\pi$, and these $gr$ integrals is linearly independent over
$\Qbar$.
\end{theorem}
\begin{proof}
The forms
\[
 \omega_k=x^k\,dx/y,\qquad0\le k\le g-1,
\]
are an algebraic basis of the regular differentials on $C_g$.  Via the
Abel--Jacobi map they give an algebraic cotangent basis on $J_g$.
For each $j$, lift the real interval $[a_j,b_j]$ to the positive branch of
$C_g$.  The vector
\[
 u_j=(M_{0,j},\ldots,M_{g-1,j})
\]
is a logarithm on $J_g$ of the algebraic divisor class of the difference
of its two endpoints.  Lemma~\ref{ii-lem-vandermonde}, applied to
$w(t)=f_g(t)^{-1/2}$ and the first $r$ coordinates, shows that the $u_j$
are linearly independent even over $\mathbb R$.  They are therefore
independent over $\mathbb Z$.  Proposition~\ref{ii-prop-endomorphisms}
and Lemma~\ref{ii-lem-logarithms} prove separation from
$\mathcal L_{\rm alg}$.

For the final assertion, a relation involving $1$, $\pi$ and the $M_{k,j}$
first has all integral coefficients equal to zero by separation.  The
remaining relation between $1$ and $\pi$ is trivial by the transcendence
of $\pi$.
\end{proof}

\begin{remark}\label{ii-rem-rg}
The restriction $r\le g$ is used exactly at the positive determinant
step.  It is a sufficient condition for this proof, not a claim that more
than $g$ intervals must yield dependence.  The construction below supplies
infinitely many independent endpoint integrals on $C_2$ by a different
argument.  In Theorem~\ref{ii-thm-intervals}, taking $r=g$ gives $g^2$
explicit independent integrals from $g$ ordered intervals on one curve.
\end{remark}

\subsection{Four independent integrals and normalized cut probabilities}

Specialize to
\[
 f(x)=x^5-x+1,\qquad
 I_k(x)=\int_0^x\frac{t^k\,dt}{\sqrt{f(t)}}\quad(k=0,1).
\]

\begin{theorem}[Every algebraic interior cut; result 71]
\label{ii-thm-four}
For every real algebraic $x$ with $0<x<1$, the four numbers
\begin{equation}\label{ii-eq-four}
 I_0(x),\quad I_1(x),\quad I_0(1),\quad I_1(1)
\end{equation}
are linearly independent over $\Qbar$, even modulo
$\mathcal L_{\rm alg}$.  Consequently,
\begin{equation}\label{ii-eq-normalized}
 F(x)=\frac{I_0(x)}{I_0(1)}
 \quad\hbox{and}\quad
 \frac{I_1(x)}{I_1(1)}
\end{equation}
are transcendental.  In fact, the ratio of any two distinct members of
\eqref{ii-eq-four} is transcendental.
\end{theorem}
\begin{proof}
Apply Theorem~\ref{ii-thm-intervals} to the two intervals $[0,x]$ and
$[x,1]$.  Its four independent integrals are
\[
 I_0(x),\ I_1(x),\ I_0(1)-I_0(x),\ I_1(1)-I_1(x).
\]
The rational change of coordinates that replaces the last two entries by
their sums with the first two is invertible.  It preserves independence
and separation from $\mathcal L_{\rm alg}$, proving the first assertion.
All four numbers are positive.  If the ratio of two distinct entries were
$c\in\Qbar$, their difference after multiplication by $c$ would be a
nontrivial algebraic linear relation, a contradiction.
\end{proof}

There is an elementary way to see the real-vector hypothesis underlying
this argument.  The ratio
\[
 \frac{I_1(x)}{I_0(x)}
\]
is the mean of $t$ on $[0,x]$ for the strictly positive weight
$f(t)^{-1/2}$.  The corresponding mean on $[x,1]$ is strictly larger.
Their two moment vectors are therefore nonproportional.  This elementary
positivity statement does not itself prove transcendence; the arithmetic
conclusion follows only after the endomorphism theorem and the analytic
subgroup theorem have been applied.

\begin{corollary}[Affine weights and quantiles]
\label{ii-cor-cdf}
Let $W(t)=a+bt\in(\Qbar\cap\mathbb R)[t]$ be nonnegative on $[0,1]$
and not identically zero.  Set
\[
 Z_W=\int_0^1\frac{W(t)\,dt}{\sqrt{f(t)}},\qquad
 F_W(x)=\frac1{Z_W}\int_0^x\frac{W(t)\,dt}{\sqrt{f(t)}}.
\]
Then $F_W(x)$ is transcendental for every algebraic $x\in(0,1)$.
For every algebraic $q\in(0,1)$, the unique solution $x\in(0,1)$ of
$F_W(x)=q$ is transcendental.
\end{corollary}
\begin{proof}
Positivity gives $Z_W>0$.  If $F_W(x)=c\in\Qbar$ at an algebraic
interior point, then
\[
 aI_0(x)+bI_1(x)-caI_0(1)-cbI_1(1)=0.
\]
The coefficients are not all zero because $(a,b)\ne(0,0)$, contrary to
Theorem~\ref{ii-thm-four}.  The distribution function is continuous and
strictly increasing: the nonzero affine polynomial $W$ cannot vanish on
an interval.  Thus each interior quantile exists uniquely.  If it were
algebraic, the preceding assertion would contradict the algebraicity of
$q$.
\end{proof}

The same argument works on $C_g$ for any nonzero polynomial weight of
degree at most $g-1$ which is nonnegative on $[0,1]$: use the $2g$
independent integrals supplied by the intervals $[0,x]$ and $[x,1]$.
This degree restriction keeps the differential in the explicitly used
space of regular differentials.  No assertion here extends this proof
automatically to arbitrary polynomial weights of larger degree.

\subsection{An infinite family on the fixed genus-two curve}

We next remove the finite-interval limitation by constructing infinitely
many independent quadratic characters.  The curve, its Jacobian, the
positive real branches, and the functions $I_0,I_1$ remain fixed.
For a prime $p$, put
\begin{equation}\label{ii-eq-Np}
 N_p=p^5-p^4+1,\qquad
 P_p=\left(\frac1p,\sqrt{\frac{N_p}{p^5}}\right),\qquad
 D_p=[P_p-\infty]\in J_2(\Qbar).
\end{equation}
The square root in $P_p$ is positive in the real embedding.

Define a strictly increasing sequence of primes recursively by
\begin{equation}\label{ii-eq-prime-sequence}
 p_1=5,\qquad
 p_n=\min\left\{p>p_{n-1}:\ p\text{ is prime},\quad
             p\nmid2869\prod_{j<n}N_{p_j}\right\}\quad(n\ge2).
\end{equation}
At each stage only finitely many primes are excluded, so the sequence is
well defined and infinite.  Its initial terms, obtained by integer
divisibility alone, are
\[
 5,7,11,13,17,23,29,31,37,47,53,59,67,71,73,\ldots.
\]
The sequence is a precise recursive definition; no assumption about
prime values of a polynomial is being made.

\begin{lemma}[A local certificate of infinite order]
\label{ii-lem-nontorsion}
If $p\ge5$ is prime and $p\nmid2869$, then the point $D_p$ in
\eqref{ii-eq-Np} has infinite order.
\end{lemma}
\begin{proof}
The polynomial discriminant is
\[
 \operatorname{disc}(x^5-x+1)=2869=19\cdot151.
\]
Since $p$ is odd and does not divide this discriminant, the smooth
projective hyperelliptic model has good reduction at $p$, as does its
Jacobian.  Work over
\[
 K_p=\mathbb Q_p(\sqrt{pN_p}).
\]
Because $N_p\equiv1\pmod p$, the polynomial $T^2-pN_p$ is Eisenstein
at $p$.  Thus $K_p/\mathbb Q_p$ is a ramified quadratic extension with
absolute ramification index $e=2$.

Near the unique point at infinity, the coordinates
\[
 u=1/x,\qquad v=y/x^3
\]
give the integral chart
\begin{equation}\label{ii-eq-infinity-chart}
 v^2=u-u^5+u^6.
\end{equation}
It is smooth at $(u,v)=(0,0)$, since the derivative of its defining
equation with respect to $u$ is a unit there.  At $P_p$ the two
coordinates are $u=p$ and $v=\sqrt{pN_p}$, so $P_p$ reduces to
$\infty$.  The Abel map therefore places $D_p$ in the kernel of
reduction on $J_2(K_p)$.  This kernel is the group of maximal-ideal
points of the formal group of the smooth Jacobian model.  Also $D_p\ne0$:
the Abel map of a positive-genus curve based at one point is injective.
Equivalently, a principal divisor $P_p-\infty$ would give a rational
function of degree one and make the curve rational.

For completeness, the relevant formal-group fact is that a commutative
formal group over the integers of a $p$-adic field has no nonzero
torsion in its maximal-ideal points when $e<p-1$
\cite[Proposition~9 and the following paragraph]{ii-ClarkXarles}.
Here is the valuation argument in the needed general dimension.
For formal coordinates $T$ one has
\[
 [p](T)=p\bigl(T+\phi(T)\bigr)+\psi(T),\qquad
 \deg\phi\ge2,\quad\deg\psi\ge p.
\]
Normalize the valuation by $v(K_p^{\times})=\mathbb Z$, so $v(p)=e$.
For a nonzero maximal-ideal tuple $T$, let $s\ge1$ be the minimum
valuation of its coordinates.  In a coordinate realizing that minimum,
the term $pT$ has valuation $e+s$, whereas the other terms have
valuation at least $e+2s$ or $ps$.  Since $e<p-1$,
$e+s<ps$, so cancellation is impossible and $[p](T)\ne0$.
Multiplication by an integer prime to $p$ is an invertible formal power
series, which excludes prime-to-$p$ torsion as well.  Any nonzero
finite-order point would then produce a nonzero point of prime order,
an impossibility.  In our case $e=2<p-1$, proving the lemma.
\end{proof}

\begin{lemma}[Independent sign changes]
\label{ii-lem-signs}
For every $r\ge1$, the classes
\[
 [f(1/p_1)],\ldots,[f(1/p_r)]
       \in\mathbb Q^{\times}/\mathbb Q^{\times2}
\]
are independent over $\mathbb F_2$.  Consequently the compositum of the
$r$ quadratic point fields has Galois group $(\mathbb Z/2\mathbb Z)^r$.
For each $j$ there is an automorphism $\sigma_j$ that changes the sign of
the ordinate of $P_{p_j}$ and fixes every $P_{p_i}$ with $i\ne j$.
On the Jacobian,
\[
 \sigma_j(D_{p_j})=-D_{p_j},\qquad
 \sigma_j(D_{p_i})=D_{p_i}\quad(i\ne j).
\]
\end{lemma}
\begin{proof}
At the new prime $p_j$, the valuation of
$f(1/p_j)=N_{p_j}/p_j^5$ is $-5$.  Its valuation on every earlier
$f(1/p_i)$ is zero by \eqref{ii-eq-prime-sequence}.  In a proposed
nonempty product of these square classes, take the largest index with
nonzero exponent.  The valuation at its corresponding prime is odd,
so the product cannot be a square.  This proves independence and the
description of the multiquadratic Galois group.

The hyperelliptic involution sends $(x,y)$ to $(x,-y)$ and fixes
$\infty$.  The divisor of $x-x(P)$ is
$P+\iota(P)-2\infty$.  Consequently it induces multiplication by
$-1$ on the Abel image based at infinity.  This proves the displayed
action on the $D_{p_i}$.
\end{proof}

\begin{theorem}[An infinite family on one curve; result 74]
\label{ii-thm-infinite}
For the sequence \eqref{ii-eq-prime-sequence}, the infinite family
\begin{equation}\label{ii-eq-infinite-family}
 I_0(1),\ I_1(1),\quad
 I_0(1/p_n),\ I_1(1/p_n)\quad(n\ge1)
\end{equation}
is linearly independent over $\Qbar$, even modulo
$\mathcal L_{\rm alg}$.  Thus adjoining $1$ and $\pi$ still leaves
a linearly independent family.  Every assertion about this infinite
family means the corresponding assertion for each finite subfamily.
\end{theorem}
\begin{proof}
Let $P_0=(0,1)$ and $P_*=(1,1)$, and use the basis
$dx/y,x\,dx/y$ of regular differentials.  Write
\[
 u_*=(I_0(1),I_1(1)),\qquad
 u_n=(I_0(1/p_n),I_1(1/p_n)).
\]
These are logarithms of the algebraic Jacobian points
\[
 Q_*=[P_*-P_0],\qquad Q_n=[P_{p_n}-P_0].
\]
We verify exactly the logarithm-vector hypothesis of
Lemma~\ref{ii-lem-logarithms}.  Suppose that for some $r$,
\begin{equation}\label{ii-eq-vector-relation}
 m_*u_*+\sum_{n=1}^r m_nu_n=0,
 \qquad m_*,m_n\in\mathbb Z.
\end{equation}
Exponentiation gives
$[m_*]Q_*+\sum_n[m_n]Q_n=0$ in $J_2(\Qbar)$.
The points $P_0,P_*$ are rational.  Put $D_0=[P_0-\infty]$; then
$Q_n=D_{p_n}-D_0$.  Applying $\sigma_j-1$ from
Lemma~\ref{ii-lem-signs} to the point relation gives
\[
 [-2m_j]D_{p_j}=0.
\]
Lemma~\ref{ii-lem-nontorsion} implies $m_j=0$ for every $j$.
Returning to the exact vector relation \eqref{ii-eq-vector-relation}
leaves $m_*u_*=0$.  Its first coordinate is strictly positive, so
$m_*=0$.  This last step does not require any assertion about whether
$Q_*$ is a torsion point: an exact relation of logarithm vectors is
stronger than the exponentiated relation.

All these logarithm vectors are therefore independent over $\mathbb Z$.
Since $\operatorname{End}(J_2)=\mathbb Z$,
Lemma~\ref{ii-lem-logarithms} proves the theorem.  The assertion with
$1$ and $\pi$ follows as in Theorem~\ref{ii-thm-intervals}.
\end{proof}

\begin{corollary}[Normalized values and pairwise ratios]
\label{ii-cor-infinite-ratios}
The family
\[
 1,\quad \frac{I_1(1)}{I_0(1)},\quad
 \frac{I_0(1/p_n)}{I_0(1)},\quad
 \frac{I_1(1/p_n)}{I_0(1)}\quad(n\ge1)
\]
is linearly independent over $\Qbar$.  The ratio of any two distinct
members of \eqref{ii-eq-infinite-family} is transcendental.
Every nonzero finite algebraic linear combination $U$ of the unnormalized
members of \eqref{ii-eq-infinite-family} is transcendental, and so are
$e^U$, $\sin U$ and $\cos U$.
\end{corollary}
\begin{proof}
For the first assertion, multiply a putative relation by the nonzero
number $I_0(1)$ and apply Theorem~\ref{ii-thm-infinite}.  An algebraic
ratio between two distinct entries would likewise give a nontrivial
linear relation.  The remaining assertions follow from separation from
$\mathcal L_{\rm alg}$ and Remark~\ref{ii-rem-limits}.
\end{proof}

\begin{remark}[Scope and provenance]\label{ii-rem-provenance}
The endomorphism theorem, the analytic subgroup theorem and the
formal-group torsion bound are classical inputs, explicitly cited above.
The displayed integral families and their independence statements are
deductions constructed from those inputs.  No priority claim is made
for their particular presentation.  Neither the finite determinant
argument nor the infinite quadratic-character construction establishes
algebraic independence of the integrals, of their normalized values, or
of their exponentials.  Also, the countably infinite independent family
on one fixed curve is not obtained by taking a limit in the
$r\le g$ theorem: its proof supplies a separate arithmetic certificate
for every finite set of endpoints.
\end{remark}



\section{Combining the two infinite families}
\label{joint-section}
The two independence constructions can be placed in one statement.
They use abelian varieties of different dimensions, so no accidental
identification of their simple factors is possible.

\begin{corollary}\label{joint-independence}
Let $\ell_p$ be the Cassini perimeters in
\eqref{lm-cassini-family}, let $p_n$ be the sequence
\eqref{ii-eq-prime-sequence}, and let $I_0,I_1$ denote the integrals
on $y^2=x^5-x+1$ in Section~\ref{ii-sec-integrals}.
Then
\[
 \{1,\pi\}\cup\{\ell_p:p\text{ an odd prime}\}
 \cup\{I_0(1),I_1(1)\}
 \cup\{I_0(1/p_n),I_1(1/p_n):n\ge1\}
\]
is linearly independent over $\Qbar$.
\end{corollary}
\begin{proof}
Consider any finite subset. Its abelian logarithms lie in a product
of pairwise nonisogenous elliptic curves $E_p$ and a power of the
absolutely simple abelian surface $J_2$. There are no nonzero
homomorphisms between $J_2$ and any $E_p$: the image of such a
homomorphism, or its dual, would contradict simplicity and the
different dimensions. In the isogeny category every connected
abelian subgroup therefore splits between these isotypic factors.
The nonzero elliptic period vectors force all elliptic factors to
occur, as in Theorem~\ref{lm-cassini-independence}. The exact
logarithm-vector independence in the proof of
Theorem~\ref{ii-thm-infinite} forces the whole power of $J_2$ to
occur. Applying the analytic subgroup theorem to a proposed linear
relation, with $\mathbb G_a$ and $\mathbb G_m$ appended for $1$ and
$2\pi i$, forces the relation to vanish on every factor. All its
coefficients are zero.
\end{proof}

This combined statement still concerns linear independence. It gives no
algebraic-independence conclusion for products, powers, or arbitrary
polynomial expressions in the displayed numbers.


\part{Transfers to other mathematical structures}

\section{Arithmetic transfer principles and applications}
\label{tf-section}

This section records the transfers numbered 1--60 in the discussion following
\cite{tf-baseline}.  Several are consequences of the same arithmetic input;
we therefore prove the common principles before giving the numbered
catalogue.  Reusing a theorem in another subject does not constitute a new
transcendence theorem.  The statements below distinguish algebraic
independence, transcendence of a single value, and effective decisions.
Throughout, algebraic means algebraic over $\mathbb Q$.

\subsection{Exponential integrals and a measure with independent moments}
\label{tf-ein-section}

Put
\[
 \Ein(z)=\int_0^z\frac{1-e^{-t}}t\,dt,
 \qquad \mathcal E=\Qbar(e^\alpha:\alpha\in\Qbar).
\]
The next argument strengthens the moment input by retaining arbitrary
algebraic exponential constants.  Its ingredients are the
Kolchin--Ostrowski theorem for antiderivatives and the
Siegel--Shidlovskii theorem for $E$-functions
\cite{tf-srinivasan,tf-fischler-rivoal}.

\begin{theorem}\label{tf-ein}
For distinct nonzero $a_1,\ldots,a_m\in\Qbar$, the numbers
$\Ein(a_1),\ldots,\Ein(a_m)$ are algebraically independent over
$\mathcal E$.
\end{theorem}
\begin{proof}
It suffices to include finitely many extra exponential constants.
Choose algebraic numbers $\beta_1,\ldots,\beta_r$ linearly independent over
$\mathbb Q$ such that every $a_j$ and every exponent occurring in these
constants is an integer combination of the $\beta_i$.  The functions
$t_i=e^{\beta_i z}$ are algebraically independent over $\mathbb C(z)$:
a polynomial relation would be a relation between exponentials with
distinct exponents and rational-function coefficients.  Let
$K=\mathbb C(z,t_1,\ldots,t_r)$, with derivation
$D=\partial_z+\sum_i\beta_i t_i\partial_{t_i}$.
Its constant field is $\mathbb C$ in its meromorphic realization.

If the primitives $y_j=\Ein(a_jz)$ were algebraically dependent over $K$,
the antiderivative theorem would give
\[
 DR=\frac{G(t_1,\ldots,t_r)}z,\qquad R\in K,
 \quad G=\sum_j c_j(1-e^{-a_jz}),\quad(c_j)\ne0,
\]
where $G$ is viewed as a Laurent polynomial in the $t_i$.
Distinct nonzero $a_j$ imply $G\ne0$.  Expand $R$ at $z=0$ over
$\mathbb C(t_1,\ldots,t_r)$.  If $R$ has pole order $d\ge1$, then
$DR$ has pole order $d+1$, because the $\partial_z$ term has strictly
larger pole order than the other terms.  If $R$ is regular, $DR$ is
regular.  A nonzero simple pole is impossible in either case.
Thus the $y_j$ are algebraically independent over $K$.

The functions $1$, $y_j$, $e^{-a_jz}$, and $e^{\beta_i z}$ form a
homogeneous linear differential system over $\Qbar(z)$ with no singularity
away from $0$.  They are $E$-functions.  Siegel--Shidlovskii at $z=1$
therefore preserves the transcendence degree $r+m$.  The exponentials
$e^{-a_j}$ and the specified extra exponentials are Laurent monomials
in the $e^{\beta_i}$.  This proves the assertion for every finite subset
of $\mathcal E$, and hence for $\mathcal E$ itself.
\end{proof}

Define a finite positive measure on $[e^{-1},1]$ by
\begin{equation}\label{tf-measure}
 d\mu(y)=\frac{1-y}{-y\log y}\,dy,
 \qquad m_n=\int y^n\,d\mu(y)
       =\Ein(n+1)-\Ein(n),\qquad \Ein(0)=0.
\end{equation}
Consequently every finite sequence $m_0,\ldots,m_N$ is algebraically
independent over $\mathcal E$.  In particular, adjoining $e$ to any such
sequence preserves algebraic independence over $\Qbar$.
A useful consequence, used repeatedly below, is
\begin{equation}\label{tf-log-rule}
 R\in\Qbar(m_0,\ldots,m_N)\setminus\Qbar,\ R>0
       \quad\Longrightarrow\quad \log R\text{ is transcendental}.
\end{equation}
Indeed, a nonconstant rational function of independent variables cannot
belong to $\mathcal E$; if $\log R$ were algebraic, $R$ would belong to
$\mathcal E$.  This includes $\log\Ein(1)$.  General results about
logarithms of $E$-function values are available in
\cite{tf-log-e-functions}; no priority claim is attached to this example.

\begin{proposition}[Quadrature coordinates]\label{tf-quadrature}
For a fixed $q\ge1$, let $x_1<\cdots<x_q$ and $w_1,\ldots,w_q>0$
be the Gaussian quadrature nodes and weights for $\mu$.
All $2q$ numbers $x_j,w_j$ are algebraically independent over
$\mathcal E$.  For $v_j=w_j/m_0$, the $2q$ numbers
\[
 m_0,x_1,\ldots,x_q,v_1,\ldots,v_{q-1}
\]
are likewise algebraically independent; the only polynomial relation
among the full normalized list is generated by $\sum_jv_j=1$.
The following further families are algebraically independent over
$\mathcal E$: all nonleading coefficients of one fixed monic orthogonal polynomial;
any initial finite list of positive Hankel determinants; and the mass
together with the corresponding finite Jacobi recurrence parameters.
\end{proposition}
\begin{proof}
Gaussian exactness gives $m_k=\sum_jw_jx_j^k$ for $0\le k<2q$.
Thus $2q$ generators contain a field of transcendence degree $2q$.
The normalized assertion is a rational change of coordinates.
The nonleading coefficients of a degree-$q$ monic polynomial generate a field over
which its $q$ roots are algebraic, proving the coefficient assertion.
Writing
$\Delta_r=\det(m_{i+j})_{0\le i,j<r}$ and $\Delta_0=1$, expansion in
$m_{2r-2}$ gives
$\Delta_r=\Delta_{r-1}m_{2r-2}+\text{earlier terms}$.
This triangular change proves independence of the $\Delta_r$.
Finally, the usual finite Jacobi recurrence parameters and $m_0$
rationally recover $m_0,\ldots,m_{2q-1}$ and conversely.
The poles and residues of the corresponding Stieltjes Pad\'e approximant
are precisely the Gaussian nodes and weights.
\end{proof}

\begin{proposition}[Errors and extremal energies]\label{tf-errors}
The $q$-node Gaussian error of a polynomial $P\in\Qbar[y]$ is zero if
$\deg P\le2q-1$, and is transcendental if $\deg P\ge2q$.
For fixed real algebraic $s$, the Christoffel extremal energies
\[
 \lambda_n(s)=\min\left\{\int P^2\,d\mu:
             P\in\mathbb R[y],\ \deg P\le n,\ P(s)=1\right\}
 \quad(n\ge0)
\]
are jointly algebraically independent, in every finite subfamily, over
$\mathcal E$.  For $q\ge2$ and integers $r\ge2$, the R\'enyi entropy
$(1-r)^{-1}\log\sum_jv_j^r$ is transcendental.
\end{proposition}
\begin{proof}
For $M\ge2q$, the nodes, weights, and errors of
$y^{2q},\ldots,y^M$ together generate a field containing all
$m_0,\ldots,m_M$.  The number of generators is $M+1$, so these
generators are algebraically independent.  Every nontrivial algebraic
linear combination of the error coordinates is transcendental.
For the second assertion, let $P_n$ be monic orthogonal polynomials and
$h_n=\int P_n^2\,d\mu$.  Then
\[
 \lambda_n(s)^{-1}-\lambda_{n-1}(s)^{-1}
          =P_n(s)^2/h_n,\qquad h_n=m_{2n}+\text{earlier terms}.
\]
Here $P_n(s)\ne0$, since its roots are transcendental.
Together with the odd moments, the $\lambda_n(s)$ therefore form a
birational coordinate system on the finite moment field.
For entropy, $v_1,\ldots,v_{q-1}$ are independent over $\mathcal E$,
with $v_q=1-\sum_{j<q}v_j$.  The positive nonconstant polynomial
$\sum_jv_j^r$ therefore cannot belong to $\mathcal E$, so its logarithm
is transcendental by the argument of \eqref{tf-log-rule}.
\end{proof}

\begin{proposition}[Polynomial ensembles and cumulants]\label{tf-ensembles}
For the probability measure
\[
 \frac1{n!\Delta_n}\prod_{i<j}(y_i-y_j)^2\prod_{i=1}^n d\mu(y_i),
\]
the partition functions $Z_n=n!\Delta_n$, $n\ge1$, are jointly algebraically
independent over $\mathcal E$, and each $\log Z_n$ is transcendental.
Every nonconstant symmetric polynomial with real algebraic coefficients
has transcendental expectation and transcendental variance.
For the normalized measure $\nu=\mu/m_0$, all its cumulants $\kappa_j$,
$j\ge1$, together
with $m_0$, are jointly algebraically independent over $\mathcal E$.
\end{proposition}
\begin{proof}
Expansion of the Vandermonde determinant expresses each expectation as
a rational function of finitely many moments.  This rational function
is nonconstant for a nonconstant symmetric polynomial $S$: specialize
the measure to exactly $n$ distinct atoms, so the ensemble is supported
on permutations of those atoms and the expectation equals the freely
varying value of $S$.  The variance is zero at such a specialization,
but is strictly positive for the continuous measure \eqref{tf-measure};
it too is a nonconstant rational function.  Independence of the moments
and \eqref{tf-log-rule} finish these assertions.
Normalized moments and cumulants are related triangularly by
$\kappa_j=m_j/m_0+\text{a polynomial in earlier normalized moments}$.
\end{proof}

\begin{corollary}[Differential equations and finite quantum dynamics]
\label{tf-dynamics}
Let $M(t)=\int e^{ty}\,d\mu(y)$ and $a_0(t)=M'(t)/M(t)$.
Neither function satisfies a nonzero algebraic differential equation
over $\Qbar(e)(t)$.  If $H$ is the $q\times q$ real symmetric Jacobi
matrix from Proposition~\ref{tf-quadrature}, $q\ge3$, its first basis
vector has no exact projective revival at any positive time under
$e^{-itH}$.
\end{corollary}
\begin{proof}
The derivatives $M^{(j)}(0)=m_j$ are algebraically independent over
$\Qbar(e)$, as are $a_0^{(j)}(0)=\kappa_{j+1}$.
Clear the denominators of a hypothetical differential equation and
remove the largest common power of $t$ from its coefficients.
Evaluation at $0$ would give a nonzero algebraic relation among these
jets.  This statement concerns the specified coefficient field;
it asserts nothing about differential algebraicity over $\mathbb C(t)$.
For a revival, positivity of all boundary spectral weights requires
$T(x_j-x_1)\in2\pi\mathbb Z$ for every $j$.  Three distinct eigenvalues
would therefore have a rational gap ratio, contradicting their
algebraic independence.  This does not exclude approximate revivals.
\end{proof}

\subsection{Polynomial beta tilts and exact cut obstructions}
\label{tf-beta-section}

Fix distinct primes $p,r$ with $\max(p,r)\ge5$, and
\begin{equation}\label{tf-beta-parameters}
 a=i/p,\quad b=j/r,\quad 0<i<p,\quad0<j<r.
\end{equation}
Let $w_0(x)=x^{a-1}(1-x)^{b-1}$, $B=B(a,b)$, and
$U(s)=B_s(a,b)$.  We use the baseline theorem that
\begin{equation}\label{tf-beta-input}
 1,B,U(s)\text{ are linearly independent over }\Qbar
       \quad(s\in\Qbar\cap(0,1)).
\end{equation}
This is the arithmetic input; the applications below do not replace it
with an assertion for arbitrary rational beta parameters.
Let $W\in(\Qbar\cap\mathbb R)[x]$ be nonzero and nonnegative on $[0,1]$.
Set $c_W=B^{-1}\int_0^1Ww_0$ and
$f_W=Ww_0/(c_WB)$.  Thus $c_W$ is positive algebraic.

\begin{lemma}[Polynomial reduction]\label{tf-beta-reduction}
For every $R\in\Qbar[x]$ there are unique $c_R\in\Qbar$ and
$q_R\in\Qbar[x]$ such that
\begin{align}
 R&=c_R+x(1-x)q_R'+[a-(a+b)x]q_R,\
 \int_0^s Rw_0\,dx&=c_RU(s)+s^a(1-s)^bq_R(s).
 \label{tf-beta-formula}
\end{align}
At algebraic $s\in(0,1)$ the normalized partial integral
$\int_0^s Rw_0/(c_WB)$ is algebraic if and only if
$c_R=q_R(s)=0$; its value is then zero.
All full polynomial moments of $f_W$ are algebraic, and are rational
when $W$ has rational coefficients.
\end{lemma}
\begin{proof}
The displayed differential operator sends $x^k$ to
$(k+a)x^k-(k+a+b)x^{k+1}$.  Its nonzero leading coefficient gives
existence and uniqueness by elimination of descending degrees, leaving
a constant.  Integration uses
$(x^a(1-x)^b)'=[a-(a+b)x]w_0$.
The criterion follows immediately from \eqref{tf-beta-input}, and the
full moments follow from
$B(a+k,b)/B(a,b)=(a)_k/(a+b)_k$.
\end{proof}

\begin{proposition}[Quantiles and conditional statistics]\label{tf-beta-stats}
For $X$ with density $f_W$, every interior quantile at an algebraic
probability is transcendental.  At every algebraic $s\in(0,1)$,
$F(s)$, $1-F(s)$, and $F(s)/(1-F(s))$ are transcendental.
If $W(s)>0$, so are the density, hazard, and reverse hazard.
For each integer $k\ge1$, both conditional moments
$\mathbb E[X^k\mid X\le s]$ and
$\mathbb E[X^k\mid X\ge s]$ are transcendental.  More precisely, each lower conditional moment
generates $\Qbar(U(s))$, and each upper conditional moment generates
$\Qbar(B-U(s))$.
\end{proposition}
\begin{proof}
Apply Lemma~\ref{tf-beta-reduction} to $W$; its partial integral is
strictly positive.  The density is a nonzero algebraic multiple of
$1/B$.  The hazard denominators are algebraic combinations of $B,U(s),1$;
\eqref{tf-beta-input} excludes an algebraic quotient.
If the lower conditional moment were algebraic $r$, applying
\eqref{tf-beta-input} to the integral of $(x^k-r)W$ would force
$r=\mathbb E[X^k]$.  This contradicts strict monotonicity of $x^k$
and positive mass on both sides of $s$.  The upper case is identical.  Formula~\eqref{tf-beta-formula}
expresses each lower conditional moment as a nonconstant fractional-linear
function of $U(s)$, and each upper one as such a function of $B-U(s)$;
inverting it gives the asserted exact fields.
Residual lifetimes follow by subtracting the algebraic cut from the
conditional first moments.  At a zero of $W$, the density and hazards
may be zero; the condition $W(s)>0$ is essential.
\end{proof}

Polynomial likelihoods, including arbitrary Bernoulli updates,
preserve this family.  Aggregated Dirichlet margins have the same
conclusion whenever their two resulting parameters meet
\eqref{tf-beta-parameters}, up to nonnegative integer shifts.
An order statistic or an iid coherent reliability system has distribution
$H(F(x))$ for a nonconstant polynomial $H$ with integer coefficients.
At an algebraic level, an algebraic value of $H(F(x))$ would imply an
algebraic value of $F(x)$; hence their interior quantiles are
transcendental.  This also applies to equal-mass grids and monotone
transport from a source distribution whose relevant cut probabilities
are algebraic.  In contrast, finite Gaussian quadrature for $f_W$ has
algebraic nodes and weights because all its full moments are algebraic.
Its finite atomic distribution cannot reproduce $F(s)$ exactly at an
algebraic interior cut.

\begin{theorem}[Optimization and multivariate polynomial tilts]
\label{tf-beta-optimization}
Let $p_0\ge1$ be an integer and $\tau\in\Qbar\cap(0,1)$.
The unique minimizer $s\in(0,1)$ of
\[
 \mathbb E\{\tau(X-s)_+^{p_0}+(1-\tau)(s-X)_+^{p_0}\}
\]
is algebraic if and only if $p_0$ is even and $\tau=1/2$.
For algebraic $s\in(0,1)$,
$W_{p_0}(\mathcal L(X),\delta_s)$ is algebraic for even $p_0$
and transcendental for odd $p_0$.
A Lloyd update from distinct algebraic centers in $(0,1)$ gives
transcendental first and last updated centers.

Moreover, tilt independent copies of a rational-parameter beta law
satisfying \eqref{tf-beta-parameters} by a nonzero nonnegative polynomial
$Q\in\mathbb Q[x_1,\ldots,x_d]$.  All polynomial joint moments remain
rational, while every marginal cumulative probability at an algebraic
interior cut is transcendental.  This applies to even-integer-$\beta$
Jacobi ensembles and their expected partial power sums.
\end{theorem}
\begin{proof}
Put $m=p_0-1$, $A=\mathbb E(X-s)_+^m$, and
$L=\mathbb E(s-X)_+^m$.  For $m=0$ these expressions mean the corresponding
side probabilities.  The optimum satisfies $\tau A=(1-\tau)L$.
If $s$ is algebraic, $A,L$ are positive transcendental partial
polynomial integrals, whereas
$M=\mathbb E(X-s)^m=A+(-1)^mL$ is algebraic.  Thus
\[
 [1-\tau+\tau(-1)^m]L=\tau M.
\]
The coefficient vanishes exactly when $p_0$ is even and $\tau=1/2$.
In that case the optimum solves the algebraic polynomial equation
$\mathbb E(X-s)^{p_0-1}=0$.  Strict convexity, or strict monotonicity of
the cumulative distribution when $p_0=1$, gives uniqueness.
The Wasserstein assertion follows by splitting $|X-s|^{p_0}$ at $s$:
for odd $p_0$ it is a full polynomial moment plus twice a nonzero
partial polynomial integral.  Algebraicity of its positive $p_0$th
root is equivalent to algebraicity of the integral.
The two outer Lloyd cells are truncated intervals with algebraic cuts,
so Proposition~\ref{tf-beta-stats} applies.  In particular no entirely
algebraic stationary codebook with at least two centers is possible;
with at least four centers, the two disjoint outer pairs force at least
two transcendental centers.

Expanding $Q$ shows that its normalizer and all joint polynomial moments
are rational.  Integrating out all but one variable leaves a nonzero
nonnegative polynomial tilt, to which the preceding proposition applies.
For an even-integer-$\beta$ Jacobi ensemble, the Vandermonde factor is
such a polynomial.  Symmetry reduces the expected partial sum
$\mathbb E\sum_i\lambda_i^m\mathbf1_{\lambda_i\le s}$, $m\ge0$, to a
positive marginal partial polynomial integral.  The conclusion does
not assert that every individual counting probability is transcendental.
\end{proof}

\subsection{Symbolic reduction, derivatives, and exact number fields}
\label{tf-symbolic-section}

For clarity we record the scope of the DBKZ input from \cite{tf-baseline}.
Let $a,b,d\in\tfrac15\mathbb Z$, $a,b>-1$, and $c>0$ algebraic.
Assume that
$(5a,5b,-5d,5(d-a-b))\bmod5$ is a permutation of
$(h,h,h,2h)$, $h\in\{1,2,3,4\}$, and that $c$ has a good reduction
with value $25$ modulo $31$, or value $7$ modulo $11$, as in the
baseline criterion.  With
\[
 g=x^a(1-x)^b(1+cx)^{-d},\qquad D=\int_0^1g\,dx,
 \qquad M=\frac1D\int_0^1xg\,dx,
\]
that criterion gives $M$ transcendental and the exact reduction
\begin{equation}\label{tf-dbkz-reduction}
 R=u+vx+\mathcal LS,\quad
 \mathcal LS=S'+\left(\frac ax-\frac b{1-x}
                         -\frac{dc}{1+cx}\right)S,
 \qquad \mathbb E_gR=u+vM.
\end{equation}
Here $R$ is rational with poles only at $0,1,-1/c,\infty$ and with
$gR$ integrable; $u,v$ are algebraic and the reduction is chosen with
vanishing endpoint term.  The baseline proves uniqueness of $u,v$
and the corresponding algebraic-primitive criterion.

\begin{proposition}\label{tf-symbolic}
All polynomial expectations in this DBKZ family generate exactly
$\Qbar(M)$.  Algebraicity of a rational expression in them is decided
by reducing it to a rational function of the indeterminate $M$:
it is algebraic precisely when that rational function is constant.
Every central moment of order $n\ge2$, and every cumulant of order
$n\ge1$, is transcendental.  An algebraic primitive of $gR$ exists
precisely when both coefficients $u,v$ in
\eqref{tf-dbkz-reduction} vanish.
\end{proposition}
\begin{proof}
Each raw moment is affine in $M$, and the first moment is $M$.
A nonconstant rational function evaluated at a transcendental number
is transcendental.  The coefficient of $M^n$ in the $n$th central
moment is $(-1)^{n-1}(n-1)$, and in the $n$th cumulant it is
$(-1)^{n-1}(n-1)!$.  These are nonzero in the stated ranges.
The primitive assertion is the two-dimensional de Rham reduction in
the baseline.  It is stronger than the definite-integral criterion
$v=0$, which permits a nonzero algebraic expectation.
\end{proof}

Two effective transfers require no new arithmetic input.  For rational $s>0$, $r>1/2$, define the generalized Brouncker constant by
\[
 Y(s,r)=4r\frac{\Gamma(A+1/2)\Gamma(B_1+1/2)}
 {\Gamma(A)\Gamma(B_1)},\qquad
 A=\frac{s+1}{4r},\quad B_1=\frac{s+2r-1}{4r}.
\]
The baseline gives this explicit finite classifier: if $2A$ or $2B_1$
is integral the value is transcendental; otherwise choose even $N$
with $NA,NB_1\in\mathbb Z$.  The value is algebraic exactly when,
for every $k$ coprime to $N$, the fractional parts $\{kA\}$ and
$\{kB_1\}$ lie in opposite open halves of $(0,1)$.  Thus equality and sign of a polynomial in $Y$ with
real algebraic coefficients are decidable.  In the algebraic cases one
uses exact algebraic arithmetic; in the transcendental cases a
nonzero polynomial cannot vanish, so certified interval approximation
eventually decides its sign.  The same criterion decides the
minimality condition $v_0/v_{-1}=s-Y(s,r)$ in
\[
 (2rn+1)v_n=2s v_{n-1}+(2rn-1)v_{n-2},
 \qquad r>1/2,\ s>0,
\]
and the corresponding sign-reversed positivity boundary
$w_0/w_{-1}=Y(s,r)-s$, with $w_{-1}>0$; the zero sequence is handled
separately.  These are decisions for the classified rational parameter
family, not for arbitrary continued fractions or recurrences.
\label{tf-decisions}

For derivative transfers put
\[
 C_h(c)=\frac{\int_0^1x^h(1-x)^h(1+cx)^{-h-1}\,dx}
                    {\int_0^1x^h(1-x)^h(1+cx)^{-h}\,dx},
 \qquad h\in\mathbb Q,\quad h>-1.
\]
We use the baseline diagonal theorem: $C_h(c)$ is transcendental for
algebraic $c>0$.  Its Riccati equation is
\begin{equation}\label{tf-riccati}
 cC_h'=(h+1)/(1+c)-C_h-hC_h^2.
\end{equation}

\begin{proposition}[Derivative and inverse transfers]\label{tf-jets}
If $h\ne0$ and $c>0$ is algebraic, then
$1,C_h(c),C_h'(c),\ldots,C_h^{(N)}(c)$ are linearly independent over
$\Qbar$ for every $N$, while their generated field has transcendence
degree one.  If $C_h(c_*)=r\in\Qbar\cap(0,1)$, then $c_*$,
$C_h'(c_*)$, and $dc_*/dr$ are transcendental.

For $F_h(c)={}_2F_1(1,h+1;2h+2;-c)$, suppose $h>-1$ has reduced
prime denominator at least $5$.  For each $n\ge1$, $F_h^{(n)}(c)$
is transcendental at all real algebraic $c>-1$, $c\ne0$, except for
at most $n$ explicitly computable algebraic points, at which it is
algebraic.  In particular, $F_{-4/5}(3)$ is transcendental but
$F_{-4/5}'(3)=-1/20$.
\end{proposition}
\begin{proof}
Successive differentiation of \eqref{tf-riccati} expresses $C_h^{(n)}$
as a polynomial in $C_h$ of degree $n+1$, with leading coefficient
$n!(-h/c)^n$.  Distinct degrees imply linear independence.
An algebraic $c_*$ would contradict the diagonal theorem.  Substitution
$C_h(c_*)=r$ in \eqref{tf-riccati} writes $C_h'(c_*)$ as a nonconstant
rational function of $c_*$ with algebraic coefficients.  It is nonzero
by strict monotonicity, so its reciprocal is transcendental as well.
For $F_h$, use
\[
 c(c+1)F_h'+[(h+1)c+2h+1]F_h=2h+1.
\]
Induction gives $F_h^{(n)}=A_nF_h+B_n$ with rational functions over
$\mathbb Q$ and $A_n=P_n(c)/[c^n(c+1)^n]$, where
$0\ne P_n\in\mathbb Q[c]$ has degree at most $n$.
Nonvanishing follows also from the homogeneous solution
$c^{-2h-1}(1+c)^h$, whose $n$th derivative has nonzero leading
asymptotic coefficient $(-1)^n(h+1)_n$.
The baseline transcendence theorem for $F_h$ gives exactly the stated
exceptional set.  Direct substitution gives the displayed exception.
\end{proof}

\begin{remark}\label{tf-jet-correction}
The exclusion $h=0$ in the first assertion is necessary: then
$cC_0'=(1+c)^{-1}-C_0$, so all jets lie in the span of $1,C_0$.
The inverse assertion does include $h=0$.
\end{remark}

Another contiguous family is
\begin{equation}\label{tf-contiguous}
 {}_2F_1(3/5+n,6/5+n;12/5+n;-1/5),\qquad n\ge0.
\end{equation}
Every value is transcendental, and these values together with $1$ span
a space of dimension exactly three over $\Qbar$.
Indeed, for $H(c)={}_2F_1(3/5,6/5;12/5;-c)$ the baseline separation
of beta and four-branch-point periods gives linear independence of
$1,H(1/5),H'(1/5)$.  The hypergeometric equation reduces all higher
jets to this span.  The differentiated Euler integrals are, up to a
nonzero algebraic factor, integrals of strictly positive rational
weights.  In their two-dimensional reduction the period coefficients
cannot both vanish: an exact term with vanishing endpoints would have
zero integral.  Thus no higher derivative is algebraic.

\begin{proposition}[The denominator-$33$ field refinement]
\label{tf-gamma33}
Let $\mathcal A_{33}$ consist of the nonidentity algebraic expressions
$Q=\Gamma(a/33)\Gamma(b/33)/[\Gamma(c/33)\Gamma(d/33)]$, with
$1\le a\le b<33$, $1\le c\le d<33$, and $a+b=c+d\ne33$,
identifying reciprocal expressions.  There are thirty such expressions
by the baseline classification.  Each quotient
$Q$ generates
\[
 K=\mathbb Q(\zeta_{33})^+(3^{1/11}),\qquad [K:\mathbb Q]=110.
\]
Of their squares, ten have degree $55$ and twenty degree $110$.
The normal closure is
$\mathbb Q(\zeta_{33},3^{1/11})$, of degree $220$, with Galois group
$\operatorname{AGL}(1,11)\times C_2$.
\end{proposition}
\begin{proof}
Write $F=\mathbb Q(\zeta_{33})^+$ and
$Q=3^{m/11}c$, where $c\in F$ and $11\nmid m$, using
Table~\ref{tf-gamma33-certificates}.  Here
\[
 t_j=2\sqrt3\sin(j\pi/33)
 =-(1+2\zeta_{33}^{11})
   \bigl((-\zeta_{33}^{17})^j-(-\zeta_{33}^{17})^{-j}\bigr)\in F.
\]
Every row follows directly from reflection and triplication,
\[
 \Gamma(j/33)\Gamma((33-j)/33)=\frac\pi{\sin(j\pi/33)},\qquad
 \frac{\Gamma(j/33)\Gamma((j+11)/33)\Gamma((j+22)/33)}
 {\Gamma(3j/33)}=2\pi\,3^{1/2-j/11}.
\]
The cyclotomic computation gives ten distinct conjugates of $c$ in
every row.  All conjugates are real, so taking eleventh powers preserves
distinctness.  Hence $\mathbb Q(Q^{11})=F\subseteq\mathbb Q(Q)$.
Since $c\in F$, the field $\mathbb Q(Q)$ contains $3^{m/11}=Q/c$;
B\'ezout's identity and $11\nmid m$ then recover $3^{1/11}$.
Thus every $Q$ generates $F(3^{1/11})$, of degree $10\cdot11=110$.

The twenty cyclotomic automorphisms applied to
$(Q^{11})^2$ give ten distinct conjugates in twenty cases and five in
the remaining ten cases, as recorded in Table~\ref{tf-gamma33-certificates}.
These counts can be checked by squaring the expressions in
Table~\ref{tf-gamma33-certificates}, raising them to the 11th power,
and reducing their displayed cyclotomic expressions modulo
$\Phi_{33}(X)=\Phi_{11}(X^3)/\Phi_{11}(X)$, with substitutions
$X\mapsto X^u$, $(u,33)=1$.  If $F_0=\mathbb Q((Q^{11})^2)$,
then $c^2\in F_0$: its conjugates over the real field $F_0$ have equal
11th powers and therefore are equal.  Hence
$\mathbb Q(Q^2)=F_0(3^{1/11})$, whose degree is $11[F_0:\mathbb Q]$.

The normal closure contains $\zeta_{11}$ and $F$, whose compositum is
$\mathbb Q(\zeta_{33})$.  Over that cyclotomic field $X^{11}-3$ is
irreducible, for the valuation of $3$ at a prime above $3$ is $2$,
coprime to $11$.  Thus the normal closure has degree $20\cdot11$.
Its group is $C_{11}\rtimes((\mathbb Z/33\mathbb Z)^\times)$,
with action through reduction modulo $11$; the order-two kernel acts
trivially, giving the asserted direct factor.
\end{proof}


\begingroup\small
\renewcommand{\arraystretch}{1.12}
\begin{longtable}{@{}llll@{}}
\caption{Exact certificates $Q=3^{m/11}c$ and square degrees for all
thirty denominator-$33$ expressions.}\label{tf-gamma33-certificates}\\
\toprule $(a,b;c,d)$ & $m$ & $c$ & $[\mathbb Q(Q^2):\mathbb Q]$\\
\midrule\endfirsthead
\toprule $(a,b;c,d)$ & $m$ & $c$ & $[\mathbb Q(Q^2):\mathbb Q]$\\
\midrule\endhead
\bottomrule\endfoot
$(1,12;3,10)$ & $-1$ & $t_{10}$ & $110$\\
$(2,13;6,9)$ & $-2$ & $t_{9}$ & $55$\\
$(3,14;8,9)$ & $-3$ & $t_{8}$ & $110$\\
$(4,15;7,12)$ & $-4$ & $t_{7}$ & $110$\\
$(5,16;6,15)$ & $-5$ & $t_{6}$ & $55$\\
$(5,18;6,17)$ & $-5$ & $\dfrac{t_{6}t_{16}}{t_{15}}$ & $110$\\
$(1,23;3,21)$ & $-1$ & $t_{12}$ & $55$\\
$(4,21;7,18)$ & $-4$ & $\dfrac{t_{7}t_{15}}{t_{12}}$ & $110$\\
$(2,24;6,20)$ & $-2$ & $t_{13}$ & $110$\\
$(3,24;8,19)$ & $-3$ & $\dfrac{t_{8}t_{14}}{t_{9}}$ & $55$\\
$(3,25;9,19)$ & $-3$ & $t_{14}$ & $110$\\
$(2,27;9,20)$ & $-2$ & $\dfrac{t_{9}t_{13}}{t_{6}}$ & $110$\\
$(4,26;12,18)$ & $-4$ & $t_{15}$ & $55$\\
$(1,30;10,21)$ & $-1$ & $\dfrac{t_{10}t_{12}}{t_{3}}$ & $110$\\
$(5,27;15,17)$ & $-5$ & $t_{16}$ & $110$\\
$(6,28;16,18)$ & $5$ & $\dfrac{t_{15}}{t_{5}t_{6}}$ & $110$\\
$(3,32;12,23)$ & $1$ & $\dfrac{1}{t_{1}}$ & $110$\\
$(7,29;15,21)$ & $4$ & $\dfrac{t_{12}}{t_{4}t_{7}}$ & $55$\\
$(6,31;13,24)$ & $2$ & $\dfrac{1}{t_{2}}$ & $110$\\
$(8,30;14,24)$ & $3$ & $\dfrac{t_{9}}{t_{3}t_{8}}$ & $110$\\
$(9,30;14,25)$ & $3$ & $\dfrac{1}{t_{3}}$ & $55$\\
$(9,31;13,27)$ & $2$ & $\dfrac{t_{6}}{t_{2}t_{9}}$ & $110$\\
$(12,29;15,26)$ & $4$ & $\dfrac{1}{t_{4}}$ & $110$\\
$(10,32;12,30)$ & $1$ & $\dfrac{t_{3}}{t_{1}t_{10}}$ & $55$\\
$(15,28;16,27)$ & $5$ & $\dfrac{1}{t_{5}}$ & $110$\\
$(17,28;18,27)$ & $5$ & $\dfrac{t_{15}}{t_{5}t_{16}}$ & $55$\\
$(18,29;21,26)$ & $4$ & $\dfrac{t_{12}}{t_{4}t_{15}}$ & $110$\\
$(19,30;24,25)$ & $3$ & $\dfrac{t_{9}}{t_{3}t_{14}}$ & $110$\\
$(20,31;24,27)$ & $2$ & $\dfrac{t_{6}}{t_{2}t_{13}}$ & $55$\\
$(21,32;23,30)$ & $1$ & $\dfrac{t_{3}}{t_{1}t_{12}}$ & $110$\\
\end{longtable}
\endgroup

Finally let $L_1=F_{1/3}(1)$ and $L_2=F_{2/3}(1)$.  Writing $g=\Gamma(1/3)$, the baseline identities are
\[
 L_1=5-\frac{40\pi^2}{3\,2^{2/3}g^3},\qquad
 L_2=\frac{21}{2}-\frac{7\sqrt3\,g^3}{6\,2^{1/3}\pi}.
\]
The algebraic independence of $\pi$ and $\Gamma(1/3)$ therefore implies
that each of $(\pi,L_1)$ and $(\pi,L_2)$ is algebraically independent.
This is a transfer of the known Chudnovsky theorem; it gives no claim
about $\Gamma(1/5)$ or Euler's constant.
\label{tf-gamma-pairs}

\subsection{Geometry, stochastic processes, and physical models}
\label{tf-physics-section}

\begin{proposition}[Closed algebraic trajectories]\label{tf-closed-curves}
Let $C$ be a smooth projective curve over the real algebraic numbers,
$\gamma$ a compact real component, and
$f:C\dashrightarrow\mathbb R^N$ an algebraic immersion without poles
on $\gamma$.  The length of $f(\gamma)$ is transcendental.  The same
holds for length in a positive algebraic metric regular on the curve.
A nonsingular compact real algebraic orbit of an algebraic vector
field has transcendental period, provided the field is everywhere
nonzero along that orbit.
\end{proposition}
\begin{proof}
The quadratic differential $q=\sum_i(df_i)^2$ has an algebraic
square-root cover, on which $\omega=\sqrt q$ is a meromorphic algebraic
differential.  Following the positive branch along the real orbit,
and going twice around if necessary, represents a positive integer
multiple of the length as a nonzero closed period.
The nonzero-period theorem for meromorphic algebraic differentials,
a consequence of the analytic subgroup theorem
\cite{lm-HuberWustholz}, makes it transcendental.
An algebraic metric changes $q$ but not this argument.
For a vector field choose the rational time form $\eta$ satisfying
$\eta(V)=1$ on the curve.  It has no pole on the orbit, and its positive
closed integral is the period.  The same theorem applies.
\end{proof}

This includes algebraic knots; for example the trigonometric algebraic
parametrization
\[
 ((2+\cos3t)\cos2t,(2+\cos3t)\sin2t,\sin3t)
\]
has squared speed $4(2+\cos3t)^2+9$ and transcendental length.
For a smooth embedded algebraic plane oval of length $L$, its complete
two-sided tube of algebraic radius $\varepsilon$ below the reach has
area $2\varepsilon L$, hence transcendental area.  This planar formula
does not automatically settle the three-dimensional tube expression
$\pi\varepsilon^2L$.
Polynomial Hamiltonian flows on regular compact level curves at
algebraic energies are included in the period statement.
Polynomial lemniscate perimeters, and hence finite polynomial escape
level perimeters, are developed more precisely in
Section~\ref{lm-section}.  In the Mandelbrot parameter setting the
finite iterates $p_0(c)=0$, $p_{n+1}(c)=p_n(c)^2+c$ are polynomials;
this concerns their algebraic level curves, not the limiting fractal
boundary.

The following direct formulas give several other transfers of the
beta input.  For the Jacobi diffusion generator
\[
 \mathcal A=\tfrac12x(1-x)\partial_x^2
  +\tfrac12[(1-a)-(2-a-b)x]\partial_x,
\]
the probability of reaching $1$ before $0$ is $U(x)/B$.
Both boundaries are accessible for $0<a,b<1$, and an algebraic
interior starting point gives a transcendental probability.
For one-dimensional steady conduction with
$\kappa(x)=x^{1-a}(1-x)^{1-b}$ and boundary values $0,1$, the solution
is $u(x)=U(x)/B$, the flux magnitude is $1/B$, and the minimum
Dirichlet energy $\int_0^1\kappa u'^2$ is $1/B$.
\label{tf-diffusion}

\begin{proposition}[Common-period resistance networks]\label{tf-networks}
In a finite connected graph, let the edge resistances be
$R_e=c_eB$, where $c_e>0$ are algebraic and $B=B(a,b)$ is as above.
Every effective resistance and every nonzero relaxation time, with
positive algebraic vertex capacities, is transcendental.
For the associated continuous-time random walk every positive-integer-order
hitting-time moment between distinct vertices is transcendental.
Its hitting probabilities and stationary probabilities are algebraic.
\end{proposition}
\begin{proof}
The conductance Laplacian is $B^{-1}L_0$ with algebraic $L_0$.
Effective resistances are positive algebraic multiples of $B$;
nonzero relaxation times have the same form.  For hitting-time moments,
the finite linear systems $Qm_n=-n m_{n-1}$, with zero target boundary
value and $Q=B^{-1}Q_0$, give $m_n=B^n a_n$ with positive algebraic
$a_n$.  The equations for hitting and stationary probabilities cancel
the common time scale.
\end{proof}

For completeness, the finite-window exponential input from
\cite{tf-baseline} says that
\begin{equation}\label{tf-window}
 J_d(c)/J_d(1)\text{ is transcendental},\qquad
 J_d(c)=\int_0^c e^{-x^{2d}}\,dx,\qquad
 d\ge1,\ c\in\Qbar\cap(0,1).
\end{equation}
It follows from the $E$-function endpoint separation theorem; a finite
mixture of $M$ algebraically specified translated and rescaled kernels
has at most $5M^2$ candidate algebraic cuts at which the normalized
cumulative probability can be algebraic, as detailed in the baseline.
This is a bound on possible exceptions, not an assertion that all
candidate points actually are exceptions.
The same input also gives a precise heat-time transfer.  For a fixed
algebraic $0<c<1$, put
\begin{equation}\label{tf-heat}
 R_c(t)=\frac{\int_0^c e^{-x^2/(4t)}\,dx}
                   {\int_0^1 e^{-x^2/(4t)}\,dx},\qquad t>0.
\end{equation}
Every solution of $R_c(t)=q\in\Qbar\cap(c,1)$ is transcendental.
Indeed, if $t$ were algebraic the rescaling $x\mapsto x/(2\sqrt t)$
and the same endpoint separation would exclude an algebraic ratio.
A solution exists because $R_c(t)$ is continuous, tends to $1$ as
$t\downarrow0$, and tends to $c$ as $t\to\infty$.
\label{tf-heat-time}

\begin{proposition}[Exact quantum probabilities]\label{tf-quantum}
On $[0,1]$ the operator
\[
 H_d=-\partial_x^2+d^2x^{4d-2}-d(2d-1)x^{2d-2},
 \quad \psi'(0)=0,\quad \psi'(1)+d\psi(1)=0,
\]
has ground state $\psi(x)=e^{-x^{2d}/2}/\sqrt{J_d(1)}$ of energy zero.
Its probability of lying in $[0,c]$ is transcendental at every algebraic
$c\in(0,1)$.  If this probability is $p$, then the two-mode state
$\sqrt p\,|10\rangle+\sqrt{1-p}\,|01\rangle$ has transcendental reduced
purity, linear entropy, and quantum Fisher information for a relative phase encoded on one mode.
No finite circuit with algebraic initial amplitudes, algebraic gates,
and algebraic measurement operators can exactly produce a specified
transcendental measurement probability.
\end{proposition}
\begin{proof}
Factorization with $A=\partial_x+dx^{2d-1}$ gives $H_d=A^*A$ with
the stated boundary conditions and $A\psi=0$.
The Born probability is \eqref{tf-window}.
The three quantities are $p^2+(1-p)^2$, $2p(1-p)$, and $4p(1-p)$,
respectively, nonconstant polynomials in a transcendental $p$.
A finite algebraic circuit retains algebraic amplitudes and hence
algebraic probabilities for computational-basis, or more generally
algebraic-operator, measurements; finite adaptivity and nonzero postselection
preserve algebraicity by sums, products, and quotients.
Approximation is not excluded.  No von Neumann or Shannon entropy
claim follows from the polynomial argument.
\end{proof}

\subsection{Numbered catalogue of transfers 1--60}
\label{tf-catalogue}

The identifiers below preserve the discussion's numbering.  ``AI'' means
algebraic independence over $\Qbar$; adjoining $e$ is included where
specified.  ``T'' means transcendental.  Statements about beta tilts use
\eqref{tf-beta-parameters}; a reference to the baseline means that the
arithmetic theorem is an input from \cite{tf-baseline}, not newly proved
by the transfer.  The later geometric developments refine item 21.

\begingroup\small
\setlength{\tabcolsep}{4pt}
\renewcommand{\arraystretch}{1.1}
\begin{longtable}{@{}L{.055\textwidth}L{.68\textwidth}L{.205\textwidth}@{}}
\toprule
No. & Audited statement and scope & Proof or input\\
\midrule\endfirsthead
\toprule No. & Audited statement and scope & Proof or input\\\midrule\endhead
\bottomrule\endfoot
1 & For fixed order, all Gaussian nodes and weights of $\mu$, together with $e$, are AI. & Proposition~\ref{tf-quadrature}\\
2 & Normalized Gaussian coordinates are AI after deleting one weight; the full list has only the sum-of-weights relation. & Proposition~\ref{tf-quadrature}\\
3 & All finite lists of positive Hankel determinants, together with $e$, are AI. & Proposition~\ref{tf-quadrature}\\
4 & The nonleading coefficients of one fixed-degree monic orthogonal polynomial, together with $e$, are AI. & Proposition~\ref{tf-quadrature}\\
5 & Finite Jacobi recurrence coordinates and the mass, together with $e$, are AI. & Proposition~\ref{tf-quadrature}\\
6 & The corresponding Stieltjes Pad\'e poles and residues, together with $e$, are AI. & Proposition~\ref{tf-quadrature}\\
7 & Every Bernoulli-updated beta posterior in the stated family has T interior algebraic-probability quantiles. & Proposition~\ref{tf-beta-stats}\\
8 & Nonnegative polynomial likelihoods preserve this quantile conclusion. & Proposition~\ref{tf-beta-stats}\\
9 & Dirichlet aggregate margins and count updates have T quantiles when the resulting two parameters meet the stated beta hypotheses. & Section~\ref{tf-beta-section}\\
10 & At algebraic interior cuts, posterior cumulative, survival, and odds values are T. & Proposition~\ref{tf-beta-stats}\\
11 & Full polynomial moments and finite polynomial predictive probabilities are algebraic; they are rational for rational polynomial weights. & Lemma~\ref{tf-beta-reduction}\\
12 & Order-statistic algebraic-level quantiles are T.  Separately, the minimum of two standard Cauchy variables has a T median. & Section~\ref{tf-beta-section}; see text below\\
13 & Iid coherent reliability systems inherit T cumulative values at algebraic cuts and T interior algebraic-level quantiles. & Section~\ref{tf-beta-section}\\
14 & Density, hazard, and reverse hazard at algebraic cuts are T if the polynomial weight is nonzero there. & Proposition~\ref{tf-beta-stats}\\
15 & Positive-order one-sided conditional moments and the corresponding mean residual lifetimes at algebraic cuts are T. & Proposition~\ref{tf-beta-stats}\\
16 & Gaussian quadrature for the beta tilts has algebraic nodes and weights, but its atomic cumulative function cannot equal the exact cumulative value at an algebraic interior cut. & Section~\ref{tf-beta-section}\\
17 & Equal-mass grids and the indicated monotone transport values are T. & Section~\ref{tf-beta-section}\\
18 & Even-integer-$\beta$ Jacobi ensembles have T expected cumulative eigenvalue fractions at algebraic interior cuts. & Theorem~\ref{tf-beta-optimization}\\
19 & Student quantiles: even integer degrees of freedom give algebraic quantiles at algebraic probabilities; odd degrees at least three give T nonmedian quantiles.  For one degree, algebraicity occurs exactly at rational probabilities. & Baseline, Student classification\\
20 & For beta shapes $\alpha,\beta>1$ rational with integer sum, shortest fixed algebraic-mass intervals have algebraic endpoints and width when both shapes are integers; otherwise endpoints and width are T with joint transcendence degree one.  In this latter case the midpoint is $1/2$ for equal shapes and T otherwise. & Baseline, HPD classification\\
21 & Every polynomial lemniscate component at a positive algebraic level has T length. & Section~\ref{lm-section}\\
22 & Finite polynomial escape-level and finite Mandelbrot-iterate level perimeters are T. & Section~\ref{tf-physics-section}\\
23 & Even-power oscillator periods and nonzero travel times between algebraic positions on a regular real branch are T. & See text below\\
24 & Internal rational quarter-phase positions and velocities of the normalized even-power oscillator are algebraic for powers $2,4,6$, and T for even powers at least $8$. & Theorem~\ref{bs-integral-classification}\\
25 & The specified Jacobi diffusion has T two-boundary exit probabilities at algebraic starting points. & Section~\ref{tf-physics-section}\\
26 & The beta-conductivity problem has T flux and minimum energy. & Section~\ref{tf-physics-section}\\
27 & An $M$-component algebraic finite-window exponential mixture has at most $5M^2$ candidate algebraic cuts with algebraic cumulative values. & Baseline, endpoint separation\\
28 & Finite-window heat fractions at algebraic target levels in $(c,1)$ can occur only at T times. & Equation~\eqref{tf-heat}\\
29 & The explicit factorized Schr\"odinger ground state has T Born probabilities at algebraic interior cuts. & Proposition~\ref{tf-quantum}\\
30 & DBKZ polynomial expectation fields equal $\Qbar(M)$; algebraicity of rational expressions is decidable by symbolic reduction. & Proposition~\ref{tf-symbolic}\\
31 & In that family, all central moments of order at least two and all positive-order cumulants are T; their field has transcendence degree one. & Proposition~\ref{tf-symbolic}\\
32 & The algebraic-primitive test is $u=v=0$, whereas algebraicity of the normalized definite integral requires only $v=0$. & Equation~\eqref{tf-dbkz-reduction}\\
33 & Equality and sign of algebraic-coefficient polynomials in the classified Brouncker constants are decidable. & Section~\ref{tf-symbolic-section}\\
34 & The stated two-parameter recurrence's minimality and associated positivity boundary are decidable at algebraic initial data in its classified rational-parameter family. & Section~\ref{tf-symbolic-section}\\
35 & For $h\ne0$, the diagonal $C_h$ derivative jets together with $1$ are linearly independent, but have transcendence degree one. & Proposition~\ref{tf-jets}\\
36 & Algebraic target values of $C_h$ have T inverse parameters, derivatives there, and inverse sensitivities. & Proposition~\ref{tf-jets}\\
37 & The $n$th derivative of the stated Log1 family has at most $n$ algebraic exceptions; $F_{-4/5}'(3)=-1/20$ is one. & Proposition~\ref{tf-jets}\\
38 & Every value in \eqref{tf-contiguous} is T; together with $1$ their linear span has dimension three. & Equation~\eqref{tf-contiguous}\\
39 & The thirty classified denominator-$33$ quotients have degree $110$; ten squares have degree $55$, twenty have degree $110$; the normal closure has degree $220$. & Proposition~\ref{tf-gamma33}\\
40 & Each of $(\pi,F_{1/3}(1))$ and $(\pi,F_{2/3}(1))$ is AI. & Section~\ref{tf-symbolic-section}\\
41 & $\log\Ein(1)$ and logarithms of positive nonconstant rational functions of the moment sequence are T. & Theorem~\ref{tf-ein}; \eqref{tf-log-rule}\\
42 & Integer-order R\'enyi entropies of Gaussian quadrature weights are T for at least two nodes and order at least two. & Proposition~\ref{tf-errors}\\
43 & Gaussian polynomial error is zero through degree $2q-1$ and T for every algebraic-coefficient polynomial of higher degree. & Proposition~\ref{tf-errors}\\
44 & At one fixed real algebraic evaluation point, all finite lists of Christoffel extremal energies, together with $e$, are AI. & Proposition~\ref{tf-errors}\\
45 & The Vandermonde-square ensemble partition functions are AI, and their logarithms are T. & Proposition~\ref{tf-ensembles}\\
46 & Every nonconstant real algebraic-coefficient symmetric polynomial observable in that ensemble has T expectation and variance. & Proposition~\ref{tf-ensembles}\\
47 & All finite lists of normalized cumulants, together with $e,m_0$, are AI. & Proposition~\ref{tf-ensembles}\\
48 & The moment-generating function and its logarithmic derivative satisfy no nonzero algebraic differential equation over $\Qbar(e)(t)$. & Corollary~\ref{tf-dynamics}\\
49 & A finite Jacobi Hamiltonian with at least three sites has no exact projective boundary-state revival at positive time. & Corollary~\ref{tf-dynamics}\\
50 & Integer-power asymmetric losses have algebraic minimizers exactly for an even power and equal side weights. & Theorem~\ref{tf-beta-optimization}\\
51 & The integer-order Wasserstein distance to an algebraic interior point mass is algebraic for even order and T for odd order. & Theorem~\ref{tf-beta-optimization}\\
52 & One Lloyd update from algebraic centers makes both outer centers T; stationary codebooks obey the stated outer-pair restrictions. & Theorem~\ref{tf-beta-optimization}\\
53 & Nonnegative rational multivariate polynomial tilts preserve rational joint moments but give T marginal algebraic-cut cumulative probabilities. & Theorem~\ref{tf-beta-optimization}\\
54 & All expected nonnegative-integer partial power sums in even-integer-$\beta$ Jacobi ensembles are T at algebraic interior cuts. & Theorem~\ref{tf-beta-optimization}\\
55 & Compact immersed closed algebraic space curves, including algebraic knots and positive algebraic metrics, have T lengths. & Proposition~\ref{tf-closed-curves}\\
56 & Thin two-sided planar tubes around embedded algebraic ovals have T areas at positive algebraic radii below the reach. & Section~\ref{tf-physics-section}\\
57 & Regular compact algebraic orbits of nonvanishing algebraic vector fields have T periods; polynomial Hamiltonian ovals are included. & Proposition~\ref{tf-closed-curves}\\
58 & Common-beta-factor networks have T effective resistances, nonzero relaxation times, and positive-integer-order intervertex hitting-time moments. & Proposition~\ref{tf-networks}\\
59 & The spatially split quantum state has T purity, linear entropy, and one-mode relative-phase quantum Fisher information. & Proposition~\ref{tf-quantum}\\
60 & Finite circuits with algebraic initial states, gates, and measurement operators cannot exactly generate a specified T measurement probability. & Proposition~\ref{tf-quantum}\\
\end{longtable}
\endgroup

Here are two details used in the catalogue.  For item 12, the median of
the minimum of two standard Cauchy variables corresponds to the
one-variable cumulative level $1-1/\sqrt2$.  Its value is
$\tan(\pi(1/2-1/\sqrt2))$, which is transcendental by
Gelfond--Schneider: algebraicity of the tangent would make the
corresponding $e^{2\pi i(1/2-1/\sqrt2)}$ algebraic, contradicting the
theorem for an algebraic irrational exponent of $-1$.
For item 23, normalize the energy equation as $\dot x^2=1-x^N$,
with $N$ a positive even integer.  For $N\ge4$, $dx/\sqrt{1-x^N}$
is holomorphic on its smooth projective hyperelliptic curve, so its
nonzero algebraic-endpoint integrals are transcendental by the Abelian
integral theorem.  For $N=2$ the integrals are nonzero algebraic
logarithms, with the same conclusion by the logarithm theorem.
Period and energy rescalings by nonzero algebraic factors preserve
transcendence.  The rational-phase division issue is substantially
stronger than this forward travel-time assertion and is treated
separately in the lemniscate division sections.

\begin{remark}[Scope of the transfers]\label{tf-scope}
The catalogue concerns exact finite constructions under explicit
hypotheses.  Transcendence is not closed under limits, multiplication
of unrelated transcendental constants, or arbitrary normalization by
another period.  The results do not prove irrationality of Euler's
constant, transcendence of $\Gamma(1/5)$, general algebraic independence
of periods, or a general statement about Shannon entropy.  Distinct
applications in the catalogue may encode the same arithmetic
obstruction rather than sixty independent new transcendence theorems.
\end{remark}
% DONE: audited transfer section and complete numbered catalogue 1--60.


\clearpage
\appendix
\section{Register of the final division and independence results}
\label{app-register}
The transfer catalogue in the preceding part records observations
1--60, and the geometric register after
Corollary~\ref{lm-cassini-ratios} records 61--69. The remaining
numbered developments are matched below to their precise statements.
The entries share methods and implications; they are not being counted
as independent solutions of unrelated open problems.

\begin{longtable}{@{}L{0.07\linewidth}L{0.56\linewidth}L{0.28\linewidth}@{}}
\toprule
No. & Content and scope & Location\\
\midrule\endfirsthead
\toprule No. & Content and scope & Location\\\midrule\endhead
70 & Rational division: elliptic reduction or a finite genus bound;
the open-arc bound includes both endpoints.
& Theorem~\ref{dd-rational}.\\[4pt]
71 & Four independent genus-two integrals at every algebraic interior
cut; normalized values and algebraic quantiles.
& Theorem~\ref{ii-thm-four}, Corollary~\ref{ii-cor-cdf}.\\[4pt]
72 & $gr$ independent moments on $r\le g$ ordered algebraic
intervals of an explicit genus-$g$ curve.
& Theorem~\ref{ii-thm-intervals}.\\[4pt]
73 & All integer-petal lemniscates: half-leaf classification,
coordinate transcendence, and exact indexed division counts.
& Corollary~\ref{bs-erdos-divisions}.\\[4pt]
74 & An explicit infinite independent family on one genus-two curve;
also jointly independent with the Cassini family.
& Theorem~\ref{ii-thm-infinite}, Corollary~\ref{joint-independence}.\\[4pt]
75 & The closed-oval bound $2g-2$ is attained in every genus;
the open-arc bound is attained in genus two.
& Theorem~\ref{dd-sharp}, Example~\ref{dd-open-sharp}.\\[4pt]
76 & Complete rational-quantile classification of
$\operatorname{Beta}(\alpha,1/2)$ for positive rational $\alpha$;
equivalent sinusoidal-spiral classification.
& Theorems~\ref{bs-beta-classification}--\ref{bs-integral-classification}.\\[4pt]
77 & All algebraic division ratios for $n=4$ and odd primes $n\ge5$,
using complete torsion packets.
& Theorem~\ref{tr-main}.\\[4pt]
78 & For general regular differentials, algebraic fractions yield
either the exact elliptic alternative or finitely many algebraic points.
& Theorem~\ref{dd-all-algebraic}.\\
\bottomrule
\end{longtable}

\section{Boundaries of the conclusions and reproducibility}
The rational-quantile classification does not imply its universal
algebraic-ratio analogue: the torsion-packet theorem covers the specified
prime cases and $n=4$, not every composite index. Finiteness does not
provide an effective enumeration or a genus bound for all algebraic ratios.
Sharpness in every genus is proved for general regular differentials,
not polynomial lemniscates.

Normalized closed traversals and branch-switching boundary routes differ.
The smooth-boundary condition for areas is essential, as the area-$2$
critical Bernoulli lemniscate shows. Elementary integrability concerns
a differential, not an individual value.

The $E$-function applications assert algebraic independence with their
specified inputs; the Abelian-integral families assert linear independence.
Neither type passes automatically to limits. No unresolved individual
Euler or Gamma constant is settled.

The standalone \LaTeX{} source embeds every section, table, proof, and
bibliography entry. No external graphics, source fragments, or bibliography
database is needed. Compile with \texttt{pdflatex} twice, or until references
stabilize. Infinite families are given by exact equations or recursions;
finite numerical lists are illustrative.

\clearpage
\begingroup\small\raggedright
\addcontentsline{toc}{section}{References}
\begin{thebibliography}{99}
\bibitem{lm-Tamiozzo}
M.~Tamiozzo,
\emph{Algebraicity of the division points of the trifolium and related topics},
arXiv:1912.00204 (2019), especially Sections~7.7 and~8.13--8.17.
\url{https://arxiv.org/abs/1912.00204}.

\bibitem{lm-HuberWustholz}
A.~Huber and G.~W\"ustholz,
\emph{Transcendence and Linear Relations of $1$-Periods},
Cambridge Tracts in Mathematics 227, Cambridge University Press, 2022.
Preprint version: arXiv:1805.10104, especially Theorem~6.2 and
Chapter~13.
\url{https://arxiv.org/abs/1805.10104}.

\bibitem{lm-Silverman}
J.~H.~Silverman,
\emph{The Arithmetic of Elliptic Curves}, second edition,
Graduate Texts in Mathematics 106, Springer, 2009,
Chapter~VII, Proposition~5.5.
\url{https://doi.org/10.1007/978-0-387-09494-6}.

\bibitem{lm-SerreTate}
J.-P.~Serre and J.~Tate,
Good reduction of abelian varieties,
\emph{Annals of Mathematics} (2) \textbf{88} (1968), 492--517.
\url{https://doi.org/10.2307/1970722}.

\bibitem{lm-Rosenlicht}
M.~Rosenlicht,
Liouville's theorem on functions with elementary integrals,
\emph{Pacific Journal of Mathematics} \textbf{24} (1968), 153--161.
\url{https://doi.org/10.2140/pjm.1968.24.153}.

\bibitem{lm-Braden}
H.~W.~Braden,
Spectral curves are transcendental,
\emph{Letters in Mathematical Physics} \textbf{111} (2021), article~9,
especially Theorem~3.1.
\url{https://doi.org/10.1007/s11005-020-01349-y}.

% lm-DONE

\bibitem{dd-Oesterle1985}
J. Oesterl\'e,
Courbes sur une vari\'et\'e ab\'elienne (d'apr\`es M. Raynaud),
\emph{Ast\'erisque} \textbf{121--122} (1985),
S\'eminaire Bourbaki 1983/84, expos\'e no.~625, 213--224.
Theorem~1 gives the finiteness of division points of a finitely generated
subgroup on a nonelliptic curve, following Raynaud and Faltings.
\url{https://www.numdam.org/item/SB_1983-1984__26__213_0/}.

\bibitem{dd-Zarhin2000}
Y. G. Zarhin,
Hyperelliptic Jacobians without complex multiplication,
\emph{Mathematical Research Letters} \textbf{7} (2000), 123--132.
\url{https://arxiv.org/abs/math/9909052}.

\bibitem{dd-Milne2013}
J. S. Milne,
\emph{Lectures on \'Etale Cohomology}, version 2.21, March 22, 2013.
Remark~3.3(c) and Theorem~3.4.
\url{https://www.jmilne.org/math/CourseNotes/LECc.pdf}.

% DD-BIB-DONE

\bibitem{bs-rosenlicht}
M.~Rosenlicht,
Generalized Jacobian varieties,
\emph{Annals of Mathematics} (2) \textbf{59} (1954), 505--530.
\url{https://doi.org/10.2307/1969715}.

\bibitem{bs-serre}
J.-P.~Serre,
\emph{Algebraic Groups and Class Fields},
Graduate Texts in Mathematics \textbf{117}, Springer, 1988,
Chapter~V.
\url{https://doi.org/10.1007/978-1-4612-1035-1}.

\bibitem{tr-boxall-grant}
J. Boxall and D. Grant,
Examples of torsion points on genus two curves,
\emph{Transactions of the American Mathematical Society}
\textbf{352} (2000), no.~10, 4533--4555.
DOI: \href{https://doi.org/10.1090/S0002-9947-00-02368-0}
{10.1090/S0002-9947-00-02368-0}.

\bibitem{tr-arul}
V. Arul,
\emph{Torsion points on Fermat quotients of the form $y^n=x^d+1$},
arXiv:1910.14251, 2019, revised 2020.
\url{https://arxiv.org/abs/1910.14251}.

\bibitem{ii-ClarkXarles}
P.~L.~Clark and X.~Xarles,
Local bounds for torsion points on abelian varieties,
\emph{Canadian Journal of Mathematics} \textbf{60} (2008), no.~3,
532--555.
See Proposition~9 and the following paragraph in the author manuscript.
\url{https://mat.uab.cat/~xarles/local.pdf}.

\bibitem{tf-baseline}
KEIO, \emph{Explicit Transcendence Criteria for Beta-Type Integrals:
Integral Quotients, Cyclotomic Fields, Statistical Quantiles, and Transfers
to the Gamma Function}, unpublished companion manuscript, 11 October 2026.

\bibitem{tf-srinivasan}
V.~R. Srinivasan, \emph{Extensions by Antiderivatives, Exponentials of
Integrals and by Iterated Logarithms}, arXiv:0811.3004, especially
Theorem~2.5. \url{https://arxiv.org/abs/0811.3004}.

\bibitem{tf-fischler-rivoal}
S.~Fischler and T.~Rivoal, \emph{Effective algebraic independence of values
of E-functions}, author manuscript, Theorem~1.
\url{https://rivoal.perso.math.cnrs.fr/articles/iaeff.pdf}.

\bibitem{tf-log-e-functions}
S.~Fischler and T.~Rivoal, \emph{Transcendence of values of logarithms of
E-functions}, arXiv:2409.18537.
\url{https://arxiv.org/abs/2409.18537}.

\end{thebibliography}
\endgroup
\end{document}
